\documentclass[twocolumn,10pt,journal]{IEEEtran}
\usepackage{amsmath,amssymb}
\usepackage{graphicx}
\usepackage{array}
\usepackage{enumitem}
\usepackage{cite}
\usepackage{url}
\usepackage[hidelinks]{hyperref}
\usepackage{orcidlink}

\begin{document}
\title{QSV2G: A Threshold Quantum Secret Sharing-Based Cross-Domain Authentication Scheme for Privacy-Preserving Vehicle-to-Grid Networks}
\author{Kumar~Prateek~\orcidlink{0000-0003-3528-3082},
Shanmukh~Datta }
\maketitle

\begin{abstract}
Cross-domain vehicle-to-grid (V2G) communication requires participant authentication and explicit authorization of charging and discharging operations. Successful authentication alone does not establish permission for an energy transfer within agreed limits and validity conditions. This paper proposes QSV2G, an architecture combining pseudonymous participant authentication with a verifiable 2-of-3 quantum secret sharing (QSS) authorization interface. The architecture distributes authorization responsibility among a Home/Registration Authority, a Visited Utility Authority, and an Independent Grid/Trust Authority. Following participant authentication and session establishment, authorization requires valid contributions from at least two distinct authorities for the same transaction context. This context binds the electric vehicle (EV) pseudonym and station identity to the session, requested operation, freshness values, authorized limits, and validity interval. The EV verifies the protected decision and authorization proof before executing the energy transfer, without receiving individual authority contributions. Routine requests use pseudonyms to limit direct disclosure of the EV's real identity. QSV2G specifies the cross-domain workflow, context-binding requirements, endpoint verification conditions, and threshold trust boundaries. The resulting contribution is an architectural and interface-level specification; concrete QSS and authorization-proof constructions, together with the adapted authentication procedure, remain necessary for end-to-end security analysis and performance evaluation.
\end{abstract}

\begin{IEEEkeywords}
Vehicle-to-grid, cross-domain authentication, quantum secret sharing, threshold authorization, pseudonymity, distributed trust.
\end{IEEEkeywords}

\section{Introduction}
\label{sec:qsv2g_introduction}

Vehicle-to-grid (V2G) communication supports coordinated energy exchange between electric vehicles (EVs) and the electricity grid. With compatible bidirectional charging equipment, an EV can receive energy for transportation and return stored energy to the grid. This capability can support demand management and participation in utility grid-service programs~\cite{doe_bidirectional}. Its operation requires communication among vehicles, charging infrastructure, and backend services. The resulting cybersecurity problem extends across organizational boundaries: the NIST cybersecurity profile for extreme fast charging considers EVs, charging infrastructure, cloud or third-party operations, and utility and building networks as interconnected domains~\cite{nist_ir8473}.

Cross-domain operation introduces a distinction between establishing participant legitimacy and granting permission for an energy transfer. Consider a registered EV requesting service from a visited charging station (CS). Successful authentication should establish the communicating participants, but it does not, by itself, determine whether that station may charge or discharge the EV, under which limits, or for how long. An authorization decision must therefore identify the permitted operation and remain bound to the participants and session for which it was issued. For example, permission to charge an EV in one session must not be interpreted as permission to discharge it in another. Routine communication should also minimize disclosure of the EV's real identity to visited infrastructure.

Quantum-related threats provide a further motivation for examining the cryptographic foundations of this communication. NIST's publication of its first three finalized post-quantum cryptography standards in 2024 establishes an important route toward addressing these threats~\cite{nist_pqc2024}. Quantum communication offers a different family of mechanisms. Within V2G research, Prateek \emph{et al.} proposed QSKA, a privacy-preserving mutual authentication scheme using superdense coding together with hash- and password-based mechanisms~\cite{prateek2024qska}. QSKA provides the authentication baseline considered in this work. Our focus is the additional authorization decision required when a pseudonymous EV requests a specific energy-transfer operation across utility domains; the authentication and enrollment procedures require adaptation to this setting.

Quantum secret sharing (QSS) provides a foundation for distributing access to secret information across multiple participants. In a threshold construction, an authorized subset can recover the protected secret, whereas an unauthorized subset cannot obtain it under the construction's security assumptions~\cite{cleve1999quantum}. This principle motivates distributing authorization responsibility rather than assigning unilateral approval capability to one authority. However, threshold secret sharing does not automatically yield a verifiable authorization protocol: the association between distributed material, context-specific contributions, and an endpoint-verifiable decision must also be defined.

To address this problem, we propose QSV2G, a cross-domain authentication and authorization architecture that combines QSKA-style pseudonymous participant authentication with a verifiable 2-of-3 QSS authorization interface. The architecture introduces three logical roles: the Home/Registration Authority (UC-A), the Visited Utility Authority (UC-B), and the Independent Grid/Trust Authority (UC-C). Following participant authentication and session establishment, authorization requires valid contributions from at least two distinct authorities for the same transaction context. The threshold policy permits progress when one authority is unavailable and the remaining two provide valid approvals. It does not provide an individual authority with veto power, and two colluding authorities lie outside the intended one-authority threshold boundary.

The authorization context binds the EV pseudonym, station identity, session identifier, requested charging or discharging operation, both participants' nonces and timestamps, the authorized limit, and the validity interval. The station returns a protected decision and a short-lived authorization proof. Before executing the operation, the EV must verify response authenticity, session and request consistency, freshness, limits, validity, the authorization proof, and an accepting decision. Individual authority contributions are not forwarded to the EV. Pseudonyms replace real identities in routine requests, although pseudonym use alone does not establish unlinkability across sessions.

\subsection{Contributions}
The contributions of this work are as follows:
\begin{enumerate}
\item \textit{Separation of authentication and operation authorization:} We define a cross-domain workflow in which participant authentication precedes an explicit authorization decision for a particular charging or discharging operation.
\item \textit{Distributed authorization architecture:} We specify three authority roles and a 2-of-3 approval policy, including distinct-authority verification, rejection when the threshold is unmet, and the policy's availability and collusion boundaries.
\item \textit{Context-bound decision and endpoint checks:} We define the authorization context, protected response fields, and EV-side acceptance conditions needed to associate permission with the intended participants, session, operation, limits, and validity.
\item \textit{Explicit realization requirements:} We identify the interfaces that a concrete implementation must supply, including pseudonym bootstrap, adapted participant authentication, verifiable QSS contributions, authorization-proof verification, and session protection.
\end{enumerate}

The present contribution is an architectural and interface-level specification. Its concrete QSS construction, authorization-token mechanism, and authentication adaptation remain to be instantiated and analyzed together. Accordingly, the specification does not establish end-to-end quantum security, unconditional security, formal verification results, or measured performance improvements. These outcomes require a complete cryptographic realization and a corresponding security and experimental evaluation.

\subsection{Organization}
Section~\ref{sec2} reviews the relevant literature. Section~\ref{sec3} defines the system and adversary models, security requirements, and assumptions. Section~\ref{sec4} presents the proposed protocol interfaces and identifies the constructions required for their realization.

\section{Related Work}
\label{sec2}

Research relevant to QSV2G spans privacy-preserving V2G authentication, cross-domain authentication and access control, post-quantum and quantum-communication-based security, and quantum secret sharing. These areas address complementary requirements: establishing participant legitimacy, protecting identity information, coordinating trust across domains, and distributing control over sensitive operations.

Gope and Sikdar~\cite{gope2019efficient} proposed a privacy-preserving authentication scheme for energy-internet-based V2G communication using lightweight primitives, including one-way hash functions. Their scheme establishes a symmetric key through which a customer accesses a service provider's services. This work illustrates the importance of combining participant authentication, key establishment, and privacy in V2G communication. For QSV2G, these functions form the basis of a protected EV--CS session. A separate design requirement is to associate an authorization decision with the particular charging or discharging operation requested within that session.

Cross-domain V2G authentication has also been investigated explicitly. Miao \emph{et al.}~\cite{miao2024threec} proposed 3C Authentication, which combines certificateless authentication with a consortium blockchain and considers both individual and aggregated verification. Kwon \emph{et al.}~\cite{kwon2024multidomain} developed a blockchain-based multi-domain V2G authentication protocol that uses hash functions and XOR operations for intra-domain authentication, while incorporating elliptic-curve cryptography and physical unclonable functions for cross-domain authentication. These studies provide relevant approaches to establishing trust when vehicles access services across administrative domains. Their use of distributed infrastructure should be distinguished from QSV2G's proposed requirement that at least two designated authorities approve the same operation context.

Authorization is itself an established topic in vehicular security. Ali \emph{et al.}~\cite{ali2023blockauth} proposed BlockAuth, a blockchain-based framework for cross-border vehicle authentication and authorization. Its smart contract evaluates service-provider access policies and incorporates separation of duty and least privilege. Although BlockAuth addresses a broader intelligent-transportation setting rather than the specific V2G energy-transfer workflow considered here, it demonstrates that authentication and policy-based authorization have already been combined in prior work. Accordingly, QSV2G's intended contribution concerns the integration of threshold approval with a pseudonymous V2G session and explicit endpoint checks; it does not originate the distinction between authentication and authorization.

Post-quantum authentication offers another relevant line of research. Dharminder and Mishra~\cite{dharminder2020lcppa} presented LCPPA, a lattice-based conditional privacy-preserving authentication framework for vehicular communication that relies on the short integer solution problem and considers service revocation. Such constructions pursue quantum resistance through computational assumptions while operating with classical cryptographic mechanisms. At the primitive level, NIST's ML-KEM standard bases key establishment on the computational difficulty of the Module Learning with Errors problem~\cite{nist2024mlkem}. Reliance on a computational assumption is therefore not, by itself, evidence that a post-quantum scheme is insecure. For the present comparison, the relevant distinction is between protecting cryptographic exchanges against quantum adversaries and distributing authority over an application-specific authorization decision.

Quantum-communication-based V2G authentication is represented directly by QSKA. Prateek \emph{et al.}~\cite{prateek2024qska} combine superdense coding, hashing, and password-based authentication and report security verification using Coq. QSKA supplies the authentication baseline considered in QSV2G. Our proposed workflow adds three logical authorization roles and requires context-consistent threshold approval after participant authentication and session establishment. This adaptation requires compatible enrollment, pseudonym availability, and authenticated session binding. Consequently, QSKA's reported verification and performance results cannot be transferred automatically to QSV2G's additional messages and authorization mechanisms.

The cryptographic foundations of distributed quantum control come from QSS research. Cleve, Gottesman, and Lo~\cite{cleve1999quantum} established quantum threshold secret sharing, in which an authorized subset can reconstruct a secret quantum state while smaller subsets obtain no information about it. Addressing malicious participation requires additional mechanisms. Ben-Or \emph{et al.}~\cite{benor2006vqss} developed verifiable quantum secret sharing and secure multiparty quantum computation with a strict honest majority, using approximate quantum error-correcting codes and quantum authentication techniques. These results provide foundations for evaluating a verifiable-QSS realization, but they do not directly define a V2G authorization token or its binding to an EV request. Likewise, collecting two classical approval messages is insufficient to instantiate a QSS construction.

Against this background, QSV2G specifies a 2-of-3 authorization interface involving a Home/Registration Authority, a Visited Utility Authority, and an Independent Grid/Trust Authority. Contributions must refer to the same pseudonym, station, session, operation, freshness values, limits, and validity interval. The EV receives a protected decision and authorization proof, rather than individual authority contributions, and must verify the response before execution. The present contribution is this architecture and its verification interfaces. A concrete QSS construction, authorization-proof mechanism, and adapted authentication procedure remain necessary to establish its security and evaluate its costs. This scope supports comparison with existing authentication, access-control, and QSS approaches without claiming demonstrated superiority or end-to-end quantum security.

\section{System Model} \label{sec3}
This section defines the system model, adversary capabilities, security and privacy requirements, assumptions, and cryptographic background used by QSV2G. The model extends the QSKA-style EV--CS authentication setting to a cross-domain V2G environment in which a roaming EV may request service from a charging station operated outside its home administrative domain. The central modeling change is that participant authentication and operation authorization are separated: EV and CS legitimacy is established through the adapted QSKA-style authentication layer, whereas permission for a particular charging or discharging transaction is distributed across three logical authorization authorities.

\subsection{Cross-Domain V2G System Model}
Consider an electric vehicle $EV_i$ registered in a home administrative domain and a visited charging station $CS_j$ operated by another utility domain. QSV2G introduces three logical authorization roles: UC-A, the Home/Registration Authority; UC-B, the Visited Utility Authority; and UC-C, an Independent Grid/Trust Authority. These roles do not replace the EV or CS authentication procedure. Instead, after successful participant authentication, they jointly determine whether the exact V2G operation requested in the current session is authorized.

Fig.~\ref{fig:qsv2g_architecture} summarizes the logical system architecture and separates participant authentication from distributed operation authorization.

\begin{figure}[!t]
\centering
\includegraphics[width=\columnwidth]{figures/qsv2g_architecture.pdf}
\caption{QSV2G logical architecture. Participant authentication precedes context-bound authorization by UC-A, UC-B, and UC-C. The station's authorization-processing block verifies contributions and applies the 2-of-3 rule; it is logic, not a fourth authority. The figure shows logical interfaces; the QSS construction and context-record transport remain to be instantiated.}
\label{fig:qsv2g_architecture}
\end{figure}

The system contains the following entities.
\begin{itemize}[leftmargin=*]
    \item \textbf{Electric Vehicle ($EV_i$):} The EV is a mobile and comparatively resource-constrained endpoint equipped with an onboard unit capable of bidirectional communication with charging infrastructure. It can operate as an energy consumer during \texttt{CHARGE} and as an energy supplier during \texttt{DISCHARGE}. During normal protocol execution, the EV is represented by a current pseudonym $PID$ rather than its real identity $I_{EV_i}$. The EV holds the credential material required by the QSKA-style participant-authentication layer and participates in fresh session establishment with the visited CS.

    \item \textbf{Visited Charging Station ($CS_j$):} The CS provides the physical interface for charging or discharging and maintains the transaction state associated with the active EV session. It is identified by $CSID$, participates in mutual participant authentication, constructs the authorization context after authentication, forwards the context to the authorization authorities, verifies returned authorization contributions, applies the threshold rule, and sends the protected authorization result to the EV. Possession of $CSID$ alone is not sufficient to establish station legitimacy.

    \item \textbf{UC-A -- Home/Registration Authority:} UC-A represents the EV's home administrative or registration domain. It maintains the home-side information required for registration, credential validation, local policy enforcement, and revocation checking. During the authorization phase it may issue one context-bound authorization contribution for the current request after its local checks succeed.

    \item \textbf{UC-B -- Visited Utility Authority:} UC-B represents the utility domain that operates or governs the visited charging infrastructure. It evaluates the current request under visited-domain policies, including operation-specific constraints, validity, limits, and revocation information available to that domain. An approving UC-B supplies its own context-bound authorization contribution.

    \item \textbf{UC-C -- Independent Grid/Trust Authority:} UC-C provides an independent authorization role separate from the home and visited utility domains. It evaluates the same request context under its applicable grid/trust policy and, when approved, supplies a third independently attributable authorization contribution.
\end{itemize}

The three authorities form a $2$-of-$3$ threshold authorization layer. Any two distinct valid contributions may satisfy the authorization threshold; one contribution alone is insufficient. Consequently, normal service can continue when one authority is unavailable and the other two provide valid approvals, while no single authority is intended to authorize a transaction independently. The threshold layer is logical authorization processing rather than a fourth utility entity. Individual QSS contributions remain within this authorization layer and are not sent to the EV.



QSV2G therefore uses two logically separate trust decisions. First, the participant-authentication layer determines whether the EV and CS are legitimate protocol participants. Second, the threshold-authorization layer determines whether enough independent authorities approve the exact operation under the current session context. Successful authentication does not by itself authorize an energy transfer, and a threshold result does not replace authentication or the EV's final verification of the protected station response.

\subsection{Adversary Model}
For classical protocol messages, QSV2G uses a Dolev--Yao-style network adversary \cite{dolev1983security}. An external adversary may eavesdrop on, intercept, delay, replay, reorder, inject, or modify messages exchanged among the EV, CS, and authorization infrastructure. The adversary may observe a pseudonym and attempt to reuse it, impersonate an EV or CS, replay an earlier authorization result, substitute one requested operation for another, alter an energy/power limit or validity interval, or move an old authorization result to a different charging station or session. The model therefore requires authentication evidence in addition to possession of public identifiers and requires authorization to be bound to the current transaction context.

The symbolic network model abstracts cryptographic operations as ideal; it does not establish their security against a quantum computer or model attacks on quantum states. A concrete realization must separately specify quantum capabilities, admissible channel operations, entanglement-distribution assumptions, and the security of the selected classical primitives against quantum adversaries. Full compromise of the EV's trusted execution logic is outside the intended endpoint-verification guarantee. A compromised CS remains an adversarial case for authorization-proof verification. Side-channel resistance and physical control of the charging hardware require separate analysis.

The distributed authorization layer introduces additional adversarial cases not present in a purely centralized trust model. A malicious or compromised authority may refuse service, remain unavailable, or return a forged, corrupted, stale, or context-mismatched authorization contribution. Such a contribution must not count toward the threshold unless it satisfies the selected verifiable-QSS verification procedure for the same authorization context. QSV2G also considers one authority to be temporarily unavailable without preventing service when the other two distinct authorities return valid contributions.

The compromise boundary is explicit. For the selected $t=2$, $n=3$ policy, one compromised authority must not be able to authorize an operation independently. If two authorities collude or are simultaneously compromised, however, they may potentially satisfy the numerical threshold. Hence, the threshold layer is intended to tolerate at most $t-1=1$ compromised authority for independent-authorization resistance, subject to the security properties of the final verifiable-QSS construction. The protocol does not claim security against an arbitrary number of compromised authorities. The final model must also define whether corruption is static, adaptive, or sequential across the lifetime of the authorization material; a numerical threshold alone does not address accumulation of compromised shares over time.

The visited CS is not granted unrestricted trust merely because it terminates the local EV connection. A malicious or impersonated station may attempt to return an unauthorized \texttt{ACCEPT} decision, but the EV is required to validate the protected session response and the corresponding authorization proof before acting on it. Because the exact $AuthToken$ construction and its verification procedure are not yet instantiated, this section states the required trust boundary rather than claiming a completed proof against a compromised CS.

Denial-of-service is treated as an availability condition rather than a confidentiality guarantee. One unavailable authority can be tolerated by the $2$-of-$3$ rule, whereas two unavailable or invalid authorities can prevent threshold authorization. No protocol can force an unavailable participant to respond, and QSV2G does not claim availability when fewer than two valid authority contributions can be obtained.

\subsection{Security and Privacy Requirements}
The security goals are derived from the cross-domain trust model and the protocol interfaces defined above. Table~\ref{secreq} lists the requirements that the QSV2G construction must satisfy. These requirements are intentionally stated independently of the still-unselected internal verifiable-QSS algorithm; the security analysis must later show whether the final instantiated construction meets them.

\begin{table*}[!t]
\caption{Security and Privacy Requirements of QSV2G}
\label{secreq}
\centering
\scriptsize
\renewcommand{\arraystretch}{1.14}
\begin{tabular}{@{}p{0.10\textwidth}p{0.20\textwidth}p{0.62\textwidth}@{}}
\hline
ID & Requirement & Definition\\
\hline
SPR1 & Participant mutual authentication & The EV and visited CS must establish the legitimacy of the communicating peer before a transaction proceeds to operation authorization.\\
SPR2 & Cross-domain authorization correctness & A V2G operation is accepted only for an authenticated transaction for which the required threshold of valid authority contributions is obtained for the same context.\\
SPR3 & Threshold enforcement & With $t=2$ and $n=3$, one authority contribution is insufficient; at least two distinct valid contributions are required for an authorization-level \texttt{ACCEPT}.\\
SPR4 & Verifiable context-bound shares & Forged, corrupted, stale, duplicate, or context-mismatched authority contributions must not be counted as valid threshold shares.\\
SPR5 & Replay and context-substitution resistance & Authorization for one session, station, operation, nonce set, time interval, or limit must not be reusable as authorization for a different context.\\
SPR6 & Fine-grained operation control & Authorization must distinguish \texttt{CHARGE} from \texttt{DISCHARGE} and bind the permitted energy/power limit and validity interval to the approved transaction.\\
SPR7 & EV identity privacy & Routine cross-domain execution must use the EV pseudonym rather than disclosing $I_{EV_i}$ to the visited CS or a passive network observer.\\
SPR8 & Protected CS-to-EV response & The final station response must provide the confidentiality/integrity/authenticity properties required by the selected session-protection mechanism and must be bound to the active session.\\
SPR9 & Single-authority compromise tolerance & Compromise of one authority must not, by itself, be sufficient to authorize a transaction under the $2$-of-$3$ threshold rule.\\
SPR10 & Single-authority unavailability tolerance & One unavailable authority must not necessarily prevent a legitimate transaction when the other two authorities provide valid contributions.\\
SPR11 & Revocation enforcement & An authority that identifies the current $PID$ as revoked must not issue an approving contribution for that request.\\
SPR12 & Controlled accountability & Normal operation remains pseudonymous, while exceptional identity disclosure, if required, must occur only through a separately authorized multi-authority procedure; its concrete cryptographic construction is outside the current core flow.\\
SPR13 & Fresh session binding & The response accepted by the EV must correspond to the active $SID$, EV/CS nonces, timestamps, requested operation, limit, validity, and authorization proof. Freshness alone is not treated as a proof of forward secrecy.\\
\hline
\end{tabular}
\end{table*}

Fresh session identifiers, nonces, and timestamps support session/context freshness but do not by themselves establish forward secrecy; that property may be claimed only if the final key-establishment mechanism provides it. Use of QSS does not automatically make the complete QSV2G protocol unconditionally secure. End-to-end security depends on the concrete QSS construction, participant-authentication procedure, key-establishment method, hash/MAC or AEAD primitive, and the assumptions under which they are analyzed.

\subsection{Assumptions}
QSV2G is developed under the following assumptions.
\begin{enumerate}[leftmargin=*]
    \item $EV_i$ and $CS_j$ complete the required enrollment/registration process before online execution and possess the credential material needed by the adapted QSKA-style authentication layer. The EV also has a valid pseudonym $PID$ available before initiating the cross-domain request; the exact bootstrap procedure that provisions this pseudonym must be fixed by the final authentication instantiation.

    \item UC-A, UC-B, and UC-C are logically distinct authorization roles. Each authority can independently evaluate the request information required by its local policy and can access the revocation state applicable to its decision. The architecture does not assume that all three authorities must approve every transaction.

    \item The selected verifiable-QSS construction realizes a $2$-of-$3$ access structure, prevents one share from independently providing authorization capability, and provides a mechanism by which forged, stale, or context-mismatched contributions can be detected before threshold counting. QSV2G does not prescribe the internal quantum share-generation, encoding, reconstruction, or commitment equations at this stage.

    \item The long-term authorization material is not exposed to the EV or treated as an ordinary reusable session credential at the CS. If the selected QSS construction requires temporary internal reconstruction, its reconstruction point, lifetime, erasure, and exposure assumptions must be defined by that concrete construction.

    \item The session-protection mechanism establishes a key $K_{EV-CS}$ that is available only to the intended authenticated EV--CS session under the assumptions of the selected key-establishment method. QSV2G requires fresh per-session binding but does not assume forward secrecy unless that method formally provides it.

    \item The hash, HMAC/encryption, or AEAD primitives used by the final implementation are selected with appropriate security parameters and are used according to their required nonce/key rules. Their standard cryptographic properties are assumed only to the extent required by the instantiated protocol.

    \item EV, CS, and authorities can perform freshness checks on nonces and timestamps. The exact clock-skew tolerance and timeout values are implementation parameters; QSV2G does not rely on timestamps alone: nonce and session freshness also require generation rules, authenticated binding, and retained transaction state.

    \item For a successful online authorization, at least two distinct authorities must be reachable and willing to return valid contributions. One unavailable authority is tolerated by design; the unavailability of two authorities may cause the request to be rejected or to time out.
\end{enumerate}

\subsection{Cryptographic Background}
QSV2G combines two different cryptographic roles: QSKA-style participant authentication and threshold QSS-based distributed operation authorization. They are complementary. The former is used to establish legitimate protocol participants and session-related secret material; the latter is introduced to distribute the authorization decision for a cross-domain V2G operation.

\subsubsection{QSKA-Style Superdense-Coding Authentication Baseline}
The QSKA baseline \cite{prateek2024qska} uses superdense coding to transfer authentication-related secret information between V2G entities. In the standard two-qubit construction, two parties share an entangled Bell pair~\cite{bennett1992dense}
\begin{equation}
\lvert\Phi^{+}\rangle=\frac{1}{\sqrt{2}}\left(\lvert00\rangle+\lvert11\rangle\right).
\label{eq:bell_pair}
\end{equation}
To encode two classical bits $(a,b)\in\{0,1\}^{2}$, the sender applies the corresponding Pauli operation to its qubit; the receiver subsequently performs the superdense-coding decoding operations and measures the two qubits to recover the encoded pair. Thus, the authentication baseline can communicate two classical bits through one transmitted qubit when the required entanglement has already been established. This accounting excludes the resources needed to distribute and maintain entanglement; dense coding alone is not an authentication or secrecy proof. QSV2G considers this QSKA-style mechanism for the participant-authentication/credential-establishment role; it does not treat superdense coding itself as the new distributed authorization mechanism.

The adapted QSKA-style authentication layer is intentionally referenced at the interface level in Section~\ref{sec4}. The exact adaptation required to make a valid pseudonym available before a roaming request, and the precise authenticated transcript binding the request fields to $Proof_{EV}$, must be fixed by the final authentication implementation. These details must not be silently inferred from the original single-domain QSKA sequence when they differ from the QSV2G cross-domain flow.

\subsubsection{Threshold QSS Authorization Abstraction}
For distributed authorization, QSV2G specifies a threshold access structure over the three authorities, motivated by QSS and verifiable-QSS constructions~\cite{cleve1999quantum,benor2006vqss}. Let
\begin{equation}
\mathcal{P}=\{A,B,C\}
\end{equation}
denote the set of authorization roles. The authorized access structure is
\begin{equation}
\Gamma=\{X\subseteq\mathcal{P}:|X|\geq2\}.
\label{eq:qss_access_structure}
\end{equation}
Equation~\eqref{eq:qss_access_structure} specifies only the intended $2$-of-$3$ authorization semantics: any two distinct authorities form an authorized set, while a single authority does not. It is not a share-generation or reconstruction algorithm.

The selected verifiable-QSS construction must eventually instantiate how long-term authorization material is prepared and distributed, how an authority produces a context-bound online contribution, how the CS-side authorization logic verifies that contribution, and how the threshold result is converted into a short-lived authorization proof. Here, $AC$ denotes the authorization-context digest and $Share_a$ denotes authority UC-$a$'s contribution; Section~\ref{sec4} defines their use. Until a construction is selected, QSV2G represents the required online verification interface abstractly as $\operatorname{Verify}(Share_a,AC)$ and leaves quantum-state preparation, reconstruction, and verification equations to that construction. This separation keeps the system and trust model precise while leaving the quantum secret-sharing primitive to be instantiated and analyzed explicitly.

\section{Proposed Protocol}
\label{sec4}

QSV2G is a privacy-preserving authentication and distributed
authorization protocol for cross-domain vehicle-to-grid (V2G)
communication. It distinguishes participant authentication
from transaction authorization, as successful authentication
alone does not permit a visited charging station (CS) to
perform unrestricted charging or discharging operations
for a roaming electric vehicle (EV). To address this,
the protocol specifies a 2-of-3 quantum-secret-sharing-assisted
authorization framework in which three utility authorities
participate in authorizing an operation under a common
session context. Pseudonyms are used during routine
communication to limit the exposure of EV identities.
The protocol covers system initialization, participant
authentication and session establishment, distributed
authorization, and protected response verification before
energy exchange. The notations used throughout QSV2G are
summarized in Table~\ref{tab:qsv2g_notation}, while
Fig.~\ref{fig:qsv2g_architecture} illustrates the
proposed protocol architecture.

\subsection{Protocol Overview and Notation}
\label{subsec:qsv2g_overview}
The participants are an electric vehicle $EV_i$, a visited charging
station $CS_j$, and three utility authorities: UC-A (home registration),
UC-B (visited utility), and UC-C (independent grid/trust). During
registration, UC-A and UC-B provision participant credentials; for an
online request, those authorities validate the respective EV and CS.
The EV and CS then establish a fresh authenticated session. Independently
of that authentication result, the three authorities evaluate an agreed
transaction context and use their reserved quantum shares only after
receiving the required local approvals. A two-authority, signed
authorization token permits the EV to verify the resulting decision
without receiving the reconstructed QSS material.

The $(2,3)$ sharing access structure applies to the authorization
resources, not automatically to all registration, consent-verification,
or session-establishment services. In the present construction UC-A is
needed online for EV credential and consent validation, and UC-B for CS
validation and proof collection. Consequently, the QSS layer supports
reconstruction by any two share holders, but uninterrupted end-to-end
operation during an arbitrary authority outage is not claimed. Quantum
state preparation and sampling checks below assume honest preparation;
they do not implement malicious-secure verifiable QSS. Table~\ref{tab:qsv2g_notation}
lists the principal notation. All hashed and authenticated multi-field
inputs use an unambiguous, domain-separated $\operatorname{Encode}$
function with fixed field ordering, typed lengths, and units.
For a candidate implementation, $H$ is SHA-256 and the HMAC/HKDF
instantiations use SHA-256; cryptographic domain-separation tags
are encoded as fixed byte strings.

\begin{table}[!t]
\caption{Main notation used in the QSV2G protocol}
\label{tab:qsv2g_notation}
\centering
\scriptsize
\renewcommand{\arraystretch}{1.07}
\begin{tabular}{@{}>{\raggedright\arraybackslash}p{0.33\columnwidth}>{\raggedright\arraybackslash}p{0.62\columnwidth}@{}}
\hline
Symbol & Description\\
\hline
$EV_i$, $CS_j$ & Vehicle and visited charging station\\
UC-A/B/C, $\mathcal P$ & Three authorities, $\mathcal P=\{A,B,C\}$\\
$PID_i$, $CSID_j$ & Registered vehicle pseudonym and station ID\\
$K_i^{reg}$, $K_j^{reg}$ & Enrollment authentication secrets\\
$K_A^{pid}$ & UC-A key for pseudonym issuance\\
$OP$, $LIM$, $VLD$ & Operation, typed limit, validity interval\\
$N_E$, $N_C$ & Fresh vehicle and station nonces\\
$TS_E$, $TS_C$ & Vehicle and station timestamps\\
$ek_i$, $dk_i$, $ct_j$ & ML-KEM keys and encapsulation ciphertext\\
$ss$, $PRK$, $SID$ & KEM secret, extracted key, session identifier\\
$K_{E\to C}$, $K_{C\to E}$ & Direction-specific AEAD traffic keys\\
$\nu$, $L$, $\mathcal Q_a^{(\nu)}$ & QSS resource ID, ternary length, share bundle\\
$\mathbf s_a^{(\nu)}$, $C_a^{(\nu)}$ & Dealer contribution and commitment\\
$\mathbf S_{auth}^{(\nu)}$ & Combined one-time classical secret\\
$M_{\mathrm{ctx}}$, $AC$ & Agreed transaction record and digest\\
$\tau_E^{AC}$ & UC-A-verifiable EV consent tag\\
$d_a(AC)$, $\alpha_a^{(\nu)}$ & Authority policy result and approval record\\
$D_{\mathrm{QSS}}$, $D_{\mathrm{auth}}$ & Reconstruction and authorization outcomes\\
$AuthToken$, $\sigma_a$ & Signed authorization token and signatures\\
$\operatorname{Encode}$, $H$ & Canonical encoding and cryptographic hash\\
$\bot$ & Failure or missing value\\
\hline
\end{tabular}
\end{table}

\subsection{System Initialization}
\label{subsec:qsv2g_initialization}
System initialization establishes the registration credentials
and quantum resources used in subsequent V2G transactions.
UC-A enrolls electric vehicles, UC-B enrolls charging stations,
and the three utility authorities prepare distributed QSS
resources for threshold authorization. The following procedure
specifies credential provisioning, Bell-state-based registration
verification, quantum-share preparation, and the state retained
by each participant upon completion.

% 1. EV Registration and Pseudonym Setup
\noindent\textbf{1) EV Registration and Pseudonym Setup:}
The electric vehicle $EV_i$ enrolls with its home utility
authority UC-A before requesting cross-domain V2G services.
The registration procedure comprises the following steps.

\textit{Step 1: Initial credential establishment.}
During secure enrollment, UC-A associates the registered
identity $I_{EV_i}$ with a secret $K_i^{reg}$, provisioned
to $EV_i$ through an authenticated and confidential channel.
The subsequent quantum-assisted exchange verifies possession
of this pre-established credential; it does not generate the
registration secret. UC-A sends a fresh challenge $r_A$
through the authenticated enrollment channel and prepares $m$
Bell pairs in the state

\begin{equation}
|\Phi^{+}\rangle_{AB}
=
\frac{|00\rangle+|11\rangle}{\sqrt{2}}.
\label{eq:ev_reg_bell}
\end{equation}

For each pair, UC-A retains qubit $B$ and transmits
qubit $A$ to $EV_i$ through the quantum channel.

\textit{Step 2: Quantum encoding.}
Upon receiving the challenge $r_A$, the EV computes
a registration authentication value

\begin{equation}
T_i =
\operatorname{Trunc}_{2m}
\left[
\operatorname{HMAC}_{K_i^{reg}}
\left(
\mathrm{REG}\parallel I_{EV_i}\parallel r_A
\right)
\right].
\label{eq:ev_reg_tag}
\end{equation}

For HMAC-SHA-256, the number of Bell pairs is constrained by
$1\leq m\leq128$, since each pair encodes two bits and the
untruncated authentication tag has 256 bits. A deployment must
select $m$ according to its authentication-security target;
smaller values are suitable only for explicitly limited tests.
The $2m$-bit authentication value is divided into
two-bit blocks $(a_k,b_k)$. For each block, the EV
encodes the corresponding Bell pair by applying

\begin{equation}
U_{a_kb_k}=Z^{a_k}X^{b_k},
\qquad a_k,b_k\in\{0,1\},
\label{eq:ev_reg_pauli}
\end{equation}

to its qubit, obtaining

\begin{equation}
|\Phi_{a_kb_k}\rangle
=
(U_{a_kb_k}\otimes I)
|\Phi^{+}\rangle.
\label{eq:ev_reg_encoded}
\end{equation}

The encoded qubits are transmitted back to UC-A.

\textit{Step 3: Quantum decoding and verification.}
UC-A performs a controlled-NOT operation, using
the returned qubit as the control, followed by a
Hadamard operation on that qubit:

\begin{equation}
|\psi_k\rangle
=
(H\otimes I)
\operatorname{CNOT}_{A\rightarrow B}
|\Phi_{a_kb_k}\rangle.
\label{eq:ev_reg_decode}
\end{equation}

Measurement in the computational basis recovers
the encoded pair $(a_k,b_k)$. UC-A concatenates the
recovered blocks to obtain $\widehat{T}_i$ and checks

\begin{equation}
\widehat{T}_i
\stackrel{?}{=}
\operatorname{Trunc}_{2m}
\left[
\operatorname{HMAC}_{K_i^{reg}}
\left(
\mathrm{REG}\parallel I_{EV_i}\parallel r_A
\right)
\right].
\label{eq:ev_reg_verify}
\end{equation}

Registration authentication succeeds only when
the recovered authentication value matches the
locally computed value for the outstanding challenge.
Both parties reject expired or already used challenges; the
Bell-state transport is not a substitute for the protected
initial provisioning of $K_i^{reg}$.

\textit{Step 4: Pseudonym provisioning.}
Following successful verification, UC-A generates
a pseudonym using a secret authority key $K_A^{pid}$,
a fresh pseudonym nonce $r_i$, and a validity
parameter $V_i$:

\begin{equation}
PID_i =
\operatorname{Trunc}_{\ell}
\left[
\operatorname{HMAC}_{K_A^{pid}}
\left(
I_{EV_i}\parallel r_i\parallel V_i
\right)
\right].
\label{eq:ev_reg_pid}
\end{equation}

UC-A stores the association between $PID_i$ and
the registered EV identity and delivers the
pseudonym to the EV through an authenticated,
confidential registration channel. The EV stores
$PID_i$ for subsequent cross-domain authentication
requests.

% 2. Charging Station Registration
\noindent\textbf{2) Charging Station Registration:}
Prior to providing V2G services, the charging station
$CS_j$ registers with its utility authority UC-B.
The registration follows the Bell-state-based authentication
procedure established for $EV_i$, with station-specific
credentials and verification parameters. The steps are
described as follows.

\textit{Step 1: Station credential establishment.}
During secure enrollment, UC-B assigns $CSID_j$ to the
charging station and provisions the shared registration secret
$K_j^{reg}$ through an authenticated and confidential channel.
UC-B retains the associated credential record. As in EV
registration, the Bell-state exchange verifies possession
of the established secret rather than generating a new one.
UC-B sends a fresh challenge $r_B$ over the protected enrollment
channel and prepares $m$ Bell
pairs according to (\ref{eq:ev_reg_bell}). For each pair,
UC-B transmits one qubit to $CS_j$ and retains its partner.

\textit{Step 2: Authentication value generation.}
Upon receiving $r_B$, the charging station computes

\begin{equation}
T_j =
\operatorname{Trunc}_{2m}
\left[
\operatorname{HMAC}_{K_j^{reg}}
\left(
\mathrm{CSREG}\parallel CSID_j\parallel r_B
\right)
\right].
\label{eq:cs_reg_auth}
\end{equation}

The resulting authentication value is divided into
$m$ two-bit blocks. For each block, $CS_j$ applies
the corresponding Pauli operation defined in
(\ref{eq:ev_reg_pauli}) to its received qubit and
returns the encoded qubits to UC-B.

\textit{Step 3: Registration verification.}
UC-B performs Bell-state decoding using the returned
qubits and their retained counterparts, following
(\ref{eq:ev_reg_decode}). The measurement outcomes
are concatenated to reconstruct $\widehat{T}_j$.
UC-B then verifies

\begin{equation}
\widehat{T}_j
\stackrel{?}{=}
\operatorname{Trunc}_{2m}
\left[
\operatorname{HMAC}_{K_j^{reg}}
\left(
\mathrm{CSREG}\parallel CSID_j\parallel r_B
\right)
\right].
\label{eq:cs_reg_verify}
\end{equation}

If the verification succeeds, UC-B confirms the
registration of $CS_j$ and retains its credential
association for subsequent authentication requests.
Otherwise, the registration attempt is rejected. A challenge
$r_B$ is accepted at most once within its enrollment window.

% 3. Distributed QSS Setup
\noindent\textbf{3) Distributed QSS Setup:}
The utility authorities UC-A, UC-B, and UC-C establish
a $(2,3)$-threshold quantum secret sharing (QSS)
arrangement using a three-qutrit encoding motivated by the
threshold construction of Cleve \emph{et al.}~\cite{cleve1999quantum}.
Here the shared data are classical ternary symbols encoded into
entangled states; the procedure does not claim to reconstruct an
arbitrary unknown quantum state.
This setup encodes one-time secret contributions into
quantum shares held separately by the authorities. Any
two designated share holders can subsequently recover
the encoded classical symbols. The procedure follows.

\textit{Authority signing credentials.}
Each UC-$a$ holds an ML-DSA-65 signing key pair;
its verification key is bound to the corresponding authority identity.
The EV and CS must obtain authentic public keys before accepting
authorization tokens. The private keys remain with their owners.

Each utility authority UC-$a$ is provisioned
with an ML-DSA-65 signing key pair

\begin{equation}
(sk_a^{sig},pk_a^{sig})
\leftarrow
\operatorname{ML\mbox{-}DSA.KeyGen}(),
\qquad a\in\mathcal{P}.
\label{eq:qsv2g_authority_signing_keys}
\end{equation}

The public verification keys are bound to
their respective authority identities through
authenticated credentials. The EV and CS
must possess a trusted means of validating
these credentials before accepting any
authority-issued authorization proof.
The private signing keys remain under the
exclusive control of their corresponding
authorities.

\textit{Step 1: Distributed secret generation.}
Let $\mathcal{P}=\{A,B,C\}$ denote the set of utility
authorities. Each authority UC-$a$, where
$a\in\mathcal{P}$, independently selects a random
ternary secret vector

\begin{equation}
\mathbf{s}_{a}^{(\nu)}
=
(s_{a,1}^{(\nu)},\ldots,s_{a,L}^{(\nu)})
\in\mathbb{Z}_3^L,
\label{eq:qss_local_secret}
\end{equation}

Here $\nu$ identifies a uniquely provisioned, single-use quantum
resource, and $L$ is the number of independent uniform ternary symbols.
In a deployment $L$ must be selected for the target secrecy level,
including resistance to quantum-accelerated guessing; small $L$ may
be used only for a labeled proof-of-concept simulation. The dealers
are assumed honest during state preparation unless a separately
validated malicious-dealer verification mechanism is added.

For each authorization resource $\nu$, UC-$a$
also computes a commitment to its secret
contribution:

\begin{equation}
\begin{aligned}
C_a^{(\nu)}
=
H\big(
&\mathrm{QSV2G\mbox{-}QSS\mbox{-}COM}
\parallel a\\
&\parallel\nu
\parallel
\operatorname{Encode}
(\mathbf{s}_a^{(\nu)})
\big).
\end{aligned}
\label{eq:qsv2g_qss_commitment}
\end{equation}

Each authority distributes its commitment
to the other authorities over authenticated
communication channels before the associated
quantum shares are used. The commitments are
retained with resource identifier $\nu$ for
subsequent reconstruction verification.
The random secret vectors must have sufficient
entropy to prevent recovery by exhaustive
search against their commitments.

The combined authorization secret is
defined componentwise as

\begin{equation}
\mathbf{S}_{auth}^{(\nu)}
=
\left(
\sum_{a\in\mathcal{P}}
\mathbf{s}_{a}^{(\nu)}
\right)\bmod 3.
\label{eq:qss_combined_secret}
\end{equation}

Each authority separately encodes and distributes its
own secret contribution using the QSS procedure below.
An authorized pair must recover the contributions from
all three authorities before computing the combined
secret in (\ref{eq:qss_combined_secret}). An individual
authority knows its own contribution, but the shares it
holds from other authorities reveal none of their
protected symbols.

\textit{Step 2: Quantum-state preparation.}
For each symbol $s_{a,k}^{(\nu)}$, UC-$a$
prepares a three-qutrit entangled state using
the encoding

\begin{equation}
\begin{aligned}
|\Psi(s)\rangle_{ABC}
=
\frac{1}{\sqrt{3}}
\sum_{r=0}^{2}
&|r\rangle_A
|r+s\rangle_B\\
&\otimes |r+2s\rangle_C,
\end{aligned}
\label{eq:qss_qutrit_encoding}
\end{equation}

where $s,r\in\mathbb{Z}_3$ and all additions
are performed modulo three. The three possible
encoded basis states are

\begin{equation}
\begin{aligned}
|\Psi(0)\rangle
&=\frac{|000\rangle+|111\rangle+|222\rangle}
{\sqrt{3}},\\
|\Psi(1)\rangle
&=\frac{|012\rangle+|120\rangle+|201\rangle}
{\sqrt{3}},\\
|\Psi(2)\rangle
&=\frac{|021\rangle+|102\rangle+|210\rangle}
{\sqrt{3}}.
\end{aligned}
\label{eq:qss_encoded_basis}
\end{equation}

Each encoded state contains three quantum
subsystems, with one subsystem designated
for each utility authority.

\textit{Step 3: Quantum-share distribution.}
Every authority distributes the three subsystems
of its encoded states according to their
corresponding authority labels. UC-$a$ retains
its designated subsystems and transmits the
remaining subsystems to the other authorities
through authenticated quantum channels.

Consequently, UC-$b$ receives a quantum-share
register from each participating authority.
Its complete share bundle for resource $\nu$
is represented as

\begin{equation}
\mathcal{Q}_{b}^{(\nu)}
=
\left\{
q_{a\rightarrow b,k}^{(\nu)}
:
a\in\mathcal{P},\
1\leq k\leq L
\right\},
\label{eq:qss_share_bundle}
\end{equation}

where $q_{a\rightarrow b,k}^{(\nu)}$
denotes the subsystem generated by UC-$a$
and assigned to UC-$b$.

For the encoding in
(\ref{eq:qss_qutrit_encoding}), the reduced
density operator of any individual subsystem
satisfies

\begin{equation}
\rho_b
=
\operatorname{Tr}_{\mathcal{P}\setminus\{b\}}
\left(
|\Psi(s)\rangle
\langle\Psi(s)|
\right)
=
\frac{I_3}{3},
\label{eq:qss_single_share_privacy}
\end{equation}

where $I_3$ is the three-dimensional identity
operator. Since this reduced state is independent of
$s$, an individual share discloses no information
about its encoded secret symbol.

For the encoded classical symbols, a pair of authorities
can recover $s$ from computational-basis measurements.
Let $x_A$, $x_B$, and $x_C$ denote the outcomes for
the respective subsystems. Equation
(\ref{eq:qss_qutrit_encoding}) gives

\begin{equation}
\begin{aligned}
x_A &= r,\\
x_B &= (r+s)\bmod 3,\\
x_C &= (r+2s)\bmod 3.
\end{aligned}
\label{eq:qss_measurement_relation}
\end{equation}

Accordingly, the reconstruction rule for each permitted
pair is

\begin{equation}
\widehat{s}=\begin{cases}
(x_B-x_A)\bmod 3, & \{A,B\},\\
2(x_C-x_A)\bmod 3, & \{A,C\},\\
(x_C-x_B)\bmod 3, & \{B,C\}.
\end{cases}
\label{eq:qss_pair_recovery}
\end{equation}

These equations specify recovery of an encoded classical
symbol, not a coherent reconstruction circuit for an
arbitrary quantum secret. Reconstruction is performed
only after the corresponding transaction is approved.

\textit{Step 4: Quantum-state consistency checks.}
Before accepting the distributed resources,
the authorities perform consistency tests on
randomly selected auxiliary encoded states.
For a qutrit, the generalized Pauli operators
are defined as

\begin{equation}
X|j\rangle=|j+1\bmod 3\rangle,
\qquad
Z|j\rangle=\omega^j|j\rangle,
\label{eq:qss_qutrit_pauli}
\end{equation}

where $\omega=e^{2\pi i/3}$.

The encoded states in
(\ref{eq:qss_qutrit_encoding}) satisfy

\begin{equation}
\begin{aligned}
(X_A\otimes X_B\otimes X_C)
|\Psi(s)\rangle
&=|\Psi(s)\rangle,\\
(Z_A\otimes Z_B\otimes Z_C)
|\Psi(s)\rangle
&=|\Psi(s)\rangle.
\end{aligned}
\label{eq:qss_state_consistency}
\end{equation}

For randomly selected auxiliary states, the authorities
measure in the corresponding generalized Pauli basis and
compare the outcomes over authenticated classical channels.
For either operator in (\ref{eq:qss_state_consistency}),
the product of the local eigenvalue outcomes must equal
unity. Tested states are discarded and cannot be used
for reconstruction. These tests check consistency with
the encoded subspace, but do not identify the encoded
symbol or establish verifiability against an arbitrary
malicious authority; a separate malicious-secure verification construction
and proof would be required for that stronger claim. The checks
also require authenticated announcement of the sampled positions
after distribution, and sampled qutrits may not be returned to the
production share pool.

The remaining unmeasured share registers are reserved as
one-time authorization resources. Following independent
approval of a common authorization context $AC$, a permitted
pair may participate in threshold reconstruction. The
transaction-specific reconstruction exchange and subsequent
proof generation belong to the distributed authorization
phase.

% 4. Initialization Completion
\noindent\textbf{4) Initialization Completion:}
After registration, $EV_i$ retains $K_i^{reg}$ and its
assigned pseudonym $PID_i$, while UC-A maintains the
associated identity and credential records. Likewise,
$CS_j$ stores $CSID_j$ and $K_j^{reg}$; UC-B retains
the corresponding registration record. Each authority
UC-$a$ maintains its assigned quantum-share bundles
$\mathcal{Q}_{a}^{(\nu)}$ indexed by authorization resource
$\nu$. Test states are discarded, while unmeasured
resources are reserved for later threshold operations.
Each resource identifier is unique and is assigned to
at most one reconstruction attempt. A share that has
been measured, rejected, or consumed is not reused.
Registration credentials and reserved quantum shares
remain under the control of their respective holders;
no long-term QSS share or complete authorization secret
is supplied to $EV_i$ or $CS_j$. The resulting state is
available to the authentication and distributed
authorization phases.

\subsection{Authentication and Session Establishment}
\label{subsec:qsv2g_authentication}

% Short opening paragraph introducing online authentication.
Following system initialization, $EV_i$ initiates an
authentication exchange with the visited charging station
$CS_j$ using its registered pseudonym $PID_i$. The
registration secrets established earlier are used to
verify the participants without repeating the enrollment
procedure. Since the EV and CS belong to different utility
domains, UC-A and UC-B participate in validating their
respective credentials. Fresh nonces and timestamps are
incorporated into the authentication messages to bind
the exchange to the current request and detect replayed
transmissions. Upon successful mutual authentication,
the EV and CS derive distinct keys for each communication
direction and a session identifier for subsequent exchanges.
The procedure consists of EV request generation,
cross-domain authentication, charging station verification,
and session key establishment.

% 1. EV Authentication Request Generation
\noindent\textbf{1) EV Authentication Request Generation:}
To request a V2G service, $EV_i$ selects an operation
$OP\in\{\mathrm{CHARGE},\mathrm{DISCHARGE}\}$ and obtains
the identifier $CSID_j$ advertised by the visited charging
station. The EV generates a fresh random nonce $N_E$
and records its current timestamp $TS_E$. Using the
registration secret $K_i^{reg}$, it computes an
authentication tag over the request parameters as

\begin{equation}
\begin{aligned}
\tau_E = \operatorname{HMAC}_{K_i^{reg}}\big(
&\mathrm{EV\_AUTH}\parallel PID_i
\parallel CSID_j\\
&\parallel OP\parallel N_E
\parallel TS_E\big).
\end{aligned}
\label{eq:ev_auth_proof}
\end{equation}

Here, the constant $\mathrm{EV\_AUTH}$ distinguishes
online authentication from the registration procedure.
Including $CSID_j$ and $OP$ binds the tag to the intended
station and requested service, while $N_E$ and $TS_E$
identify the current authentication attempt. All fields
are encoded unambiguously before the HMAC is evaluated.

The EV then constructs the authentication request

\begin{equation}
\begin{aligned}
Req_{EV} = (&PID_i, CSID_j, OP,\\
            &N_E, TS_E, \tau_E),
\end{aligned}
\label{eq:ev_auth_request}
\end{equation}

and transmits $Req_{EV}$ to $CS_j$. Since the station
does not possess $K_i^{reg}$, it cannot independently
verify $\tau_E$. Instead, the request is submitted
for validation by UC-A through the cross-domain
authentication procedure described next.

% 2. Cross-Domain EV Authentication
\noindent\textbf{2) Cross-Domain EV Authentication:}
Upon receiving $Req_{EV}$, the visited charging station
$CS_j$ initiates verification through UC-B, which forwards
the authentication request to the EV's home authority,
UC-A. The exchange between utility authorities takes
place over authenticated and confidential communication
channels. UC-A verifies the EV using the registration
information established during system initialization.
The procedure is performed as follows.

\textit{Step 1: Authentication request forwarding.}
The charging station forwards the received request
(\ref{eq:ev_auth_request}) to UC-B together with its
station identifier $CSID_j$. UC-B relays the request
to UC-A without modifying the authentication fields.
The identity of the home authority is resolved through
the registered pseudonym-routing information. The
vehicle's real identity is not included in the forwarded
request.

\textit{Step 2: Credential and freshness verification.}
On receiving the request, UC-A locates the registration
record associated with $PID_i$ and retrieves the
corresponding secret $K_i^{reg}$. The authority first
checks whether the pseudonym remains valid and whether
the registration record has been revoked.

To reject stale or previously processed requests,
UC-A verifies the timestamp and nonce according to

\begin{equation}
\begin{aligned}
\operatorname{Fresh}_{E}
=
\bigl(&|T_A-TS_E|\leq\Delta_T\bigr)
\\
&\land
\bigl(N_E\notin\mathcal{N}_i\bigr),
\end{aligned}
\label{eq:ev_auth_freshness}
\end{equation}

where $T_A$ denotes the verification time at UC-A,
$\Delta_T$ is the maximum permitted clock deviation
and transmission delay, and $\mathcal{N}_i$ contains
the previously accepted nonces associated with
$PID_i$ within the replay-protection window.
An invalid timestamp or repeated nonce causes
the authentication attempt to be rejected.

\textit{Step 3: Authentication proof verification.}
After the freshness checks, UC-A independently
computes the expected authentication tag using
the stored registration secret:

\begin{equation}
\begin{aligned}
\tau_E^{*}
=
\operatorname{HMAC}_{K_i^{reg}}\big(
&\mathrm{EV\_AUTH}\parallel PID_i
\parallel CSID_j\\
&\parallel OP\parallel N_E
\parallel TS_E\big).
\end{aligned}
\label{eq:ev_auth_recomputed}
\end{equation}

The received tag $\tau_E$ is then compared with
$\tau_E^{*}$ using a constant-time comparison.
The EV is accepted only if its registration record
is valid, the freshness requirements are satisfied,
and the authentication tags match:

\begin{equation}
\begin{aligned}
\operatorname{Auth}_{EV}
={}&\operatorname{Valid}(PID_i)
\\
&\land\operatorname{Fresh}_{E}
\\
&\land
(\tau_E\stackrel{?}{=}\tau_E^{*}).
\end{aligned}
\label{eq:ev_auth_decision}
\end{equation}

Here, $\operatorname{Valid}(PID_i)$ denotes
the successful registration-status and
pseudonym-validity checks. Following successful
verification, UC-A atomically records $N_E$
in the replay-protection state for $PID_i$,
preventing the same request from being accepted
again.

\textit{Step 4: Authentication result forwarding.}
UC-A constructs a session-bound authentication
result containing the verified pseudonym,
station identifier, requested operation, and
EV nonce:

\begin{equation}
\begin{aligned}
R_A = (&PID_i,CSID_j,OP,N_E,\\
       &TS_E,\operatorname{Auth}_{EV}).
\end{aligned}
\label{eq:ev_auth_result}
\end{equation}

UC-A transmits $R_A$ to UC-B through their
authenticated channel. UC-B forwards the result
to $CS_j$ through its protected station
communication channel. The receiving parties
accept the result only from the authenticated
authority and associate it with the corresponding
pending request.

A successful result establishes that UC-A has
verified the EV's registered credentials and
the freshness of its request. It does not,
however, authenticate the charging station
to the EV or authorize the requested energy
operation. Station authentication and subsequent
session establishment are performed in the
following steps.

% 3. Charging Station Authentication
\noindent\textbf{3) Charging Station Authentication:}
After receiving a successful EV authentication result
$R_A$, the charging station $CS_j$ must establish its
legitimacy to the requesting vehicle. Since the EV
does not possess the station's registration secret
$K_j^{reg}$, station authentication is performed
through UC-B, with UC-A providing an EV-verifiable
confirmation. The procedure consists of the
following steps.

\textit{Step 1: Station authentication proof generation.}
The charging station generates a fresh random nonce
$N_C$ and records its current timestamp $TS_C$.
Using its previously established registration
secret $K_j^{reg}$, $CS_j$ computes

\begin{equation}
\begin{aligned}
\tau_C =
\operatorname{HMAC}_{K_j^{reg}}\big(
&\mathrm{CS\_AUTH}\parallel CSID_j
\parallel PID_i\\
&\parallel OP\parallel N_E
\parallel N_C\parallel TS_C
\big).
\end{aligned}
\label{eq:cs_auth_proof}
\end{equation}

The value $\mathrm{CS\_AUTH}$ provides domain
separation from the registration and EV
authentication procedures. The inclusion of
$PID_i$, $N_E$, and $OP$ binds the station's
proof to the EV request already validated by
UC-A. The charging station sends

\begin{equation}
M_C =
(CSID_j,PID_i,OP,N_E,N_C,TS_C,\tau_C)
\label{eq:cs_auth_message}
\end{equation}

to UC-B through their authenticated
communication channel.

\textit{Step 2: Station credential verification.}
Upon receiving $M_C$, UC-B retrieves the
registration secret associated with $CSID_j$
and checks the station's registration and
revocation status. The authority also verifies
that the request fields correspond to a
successful and pending EV authentication
result $R_A$.

To prevent the acceptance of stale or repeated
station authentication messages, UC-B checks

\begin{equation}
\begin{aligned}
\operatorname{Fresh}_{C}
=
\bigl(&|T_B-TS_C|\leq\Delta_T\bigr)
\\
&\land
\bigl(N_C\notin\mathcal{N}_j\bigr),
\end{aligned}
\label{eq:cs_auth_freshness}
\end{equation}

where $T_B$ denotes the current time at UC-B,
and $\mathcal{N}_j$ maintains the accepted
station nonces within the replay-protection
window.

UC-B independently recomputes the expected
authentication tag as

\begin{equation}
\begin{aligned}
\tau_C^{*} =
\operatorname{HMAC}_{K_j^{reg}}\big(
&\mathrm{CS\_AUTH}\parallel CSID_j
\parallel PID_i\\
&\parallel OP\parallel N_E
\parallel N_C\parallel TS_C
\big).
\end{aligned}
\label{eq:cs_auth_expected}
\end{equation}

The charging station is authenticated only
when its registration record is valid,
the associated EV authentication result is
successful, the freshness checks pass,
and the authentication tags match:

\begin{equation}
\begin{aligned}
\operatorname{Auth}_{CS}
={}&\operatorname{Valid}(CSID_j)
\\
&\land\operatorname{Bind}(M_C,R_A)
\\
&\land\operatorname{Fresh}_{C}
\\
&\land(\tau_C\stackrel{?}{=}\tau_C^{*}).
\end{aligned}
\label{eq:cs_auth_decision}
\end{equation}

Here, $\operatorname{Bind}(M_C,R_A)$ verifies
the agreement of $PID_i$, $CSID_j$, $OP$,
and $N_E$ across the two messages. UC-B
atomically records $N_C$ upon successful
verification.

\textit{Step 3: Cross-domain authentication confirmation.}
Following successful verification, UC-B sends
the authenticated station-validation result

\begin{equation}
\begin{aligned}
R_B = (&PID_i,CSID_j,OP,N_E,\\
       &N_C,TS_C,\operatorname{Auth}_{CS})
\end{aligned}
\label{eq:cs_auth_result}
\end{equation}

to UC-A through the protected inter-authority
channel. UC-A checks that $R_B$ corresponds
to the EV request authenticated earlier.
It then generates an EV-verifiable confirmation

\begin{equation}
\begin{aligned}
\tau_A =
\operatorname{HMAC}_{K_i^{reg}}\big(
&\mathrm{CS\_CONF}\parallel PID_i
\parallel CSID_j\\
&\parallel OP\parallel N_E
\parallel N_C\parallel TS_C\\
&\parallel\operatorname{Auth}_{CS}
\big).
\end{aligned}
\label{eq:cs_auth_confirmation}
\end{equation}

UC-A returns the confirmation together with
the associated authentication fields through
UC-B and $CS_j$ to the EV.

\textit{Step 4: EV-side confirmation verification.}
On receiving the confirmation, $EV_i$
recomputes the expected value of $\tau_A$
using its registration secret $K_i^{reg}$.
The EV checks the authentication tag,
the validity of $TS_C$, the received station
identifier, and the agreement of $PID_i$,
$OP$, and $N_E$ with its outstanding request.
It also requires
$\operatorname{Auth}_{CS}=1$.

If these checks succeed, the EV accepts
UC-A's confirmation that UC-B has verified
the registered station for the current
authentication exchange. The nonce $N_C$
is retained for session establishment.
Proof of possession of the resulting session
key by the communicating endpoint is
completed through the subsequent key
confirmation procedure. No charging or
discharging operation is authorized at
this stage.

% 4. Session Key Generation and Establishment
\noindent\textbf{4) Session Key Generation and Establishment:}
Following the successful authentication of $EV_i$ and
$CS_j$, the participants establish a fresh session key
using the module-lattice-based key encapsulation mechanism
ML-KEM-768. The authenticated registration credentials
are used to bind the key-establishment exchange to the
previously verified participants. The procedure is
described below.

\textit{Step 1: EV key generation.}
The EV generates a fresh ML-KEM key pair for the current
authentication exchange:

\begin{equation}
(ek_i,dk_i)
\leftarrow
\operatorname{ML\mbox{-}KEM.KeyGen}().
\label{eq:qsv2g_kem_keygen}
\end{equation}

Here, $ek_i$ is the encapsulation key and $dk_i$ is
the corresponding decapsulation key. To authenticate
the encapsulation key, $EV_i$ computes

\begin{equation}
\begin{aligned}
M_E^K=\operatorname{Encode}(&PID_i,CSID_j,OP,N_E,\\
&N_C,TS_E,TS_C,ek_i),\\
\tau_E^K=\operatorname{HMAC}_{K_i^{reg}}(&\mathrm{KEM\mbox{-}EV}
\parallel H(M_E^K)).
\end{aligned}
\label{eq:qsv2g_kem_ev_proof}
\end{equation}

After validating UC-A's station confirmation and $N_C$,
the EV sends $(ek_i,\tau_E^{K})$ to $CS_j$. The station
forwards these values through UC-B to UC-A.
UC-A verifies the tag using the registered EV
credential and confirms that the values belong
to the pending authentication exchange.
The validated encapsulation key is subsequently
made available to $CS_j$ through UC-B's
authenticated channel.

\textit{Step 2: Session-secret encapsulation.}
After receiving the validated encapsulation key,
$CS_j$ invokes the ML-KEM encapsulation algorithm:

\begin{equation}
(ct_j,ss_j)
\leftarrow
\operatorname{ML\mbox{-}KEM.Encaps}(ek_i),
\label{eq:qsv2g_kem_encaps}
\end{equation}

where $ct_j$ denotes the encapsulation ciphertext
and $ss_j$ is the newly generated shared secret.

The station constructs a canonical session transcript

\begin{equation}
\begin{aligned}
\mathcal{T} =
\operatorname{Encode}(&PID_i,CSID_j,OP,N_E,N_C,\\
                     &TS_E,TS_C,ek_i,ct_j),
\end{aligned}
\label{eq:qsv2g_session_transcript}
\end{equation}

where $\operatorname{Encode}$ denotes an unambiguous,
length-delimited representation of the included
fields. The station authenticates the transcript
with its registration key:

\begin{equation}
\tau_C^{K}
=
\operatorname{HMAC}_{K_j^{reg}}
\left(
\mathrm{KEM\mbox{-}CS}\parallel H(\mathcal{T})
\right).
\label{eq:qsv2g_kem_cs_proof}
\end{equation}

The ciphertext and authentication tag are sent
to UC-B for verification.

\textit{Step 3: Authenticated ciphertext delivery.}
UC-B verifies $\tau_C^{K}$ against the registered
station credential and confirms that the transcript
matches the previously authenticated request.
Following successful verification, UC-B forwards
the validated transcript and station-authentication
result to UC-A over the protected inter-authority
channel.

UC-A checks the association between the transcript
and the authenticated EV request. It then generates
an EV-verifiable confirmation:

\begin{equation}
\tau_A^{K}
=
\operatorname{HMAC}_{K_i^{reg}}
\left(
\mathrm{KEM\mbox{-}CONF}
\parallel H(\mathcal{T})
\parallel \mathrm{ACCEPT}
\right).
\label{eq:qsv2g_kem_authority_conf}
\end{equation}

UC-A returns $\tau_A^{K}$ through UC-B and
$CS_j$ to the EV. The station also forwards
$ct_j$ and the transcript fields required
for independent verification.

Before decapsulation, $EV_i$ verifies
$\tau_A^{K}$, checks the agreement of $PID_i$,
$CSID_j$, $OP$, and $N_E$ with its pending
request, and validates the received station
nonce and timestamp.

If these checks succeed, the EV decapsulates
the received ciphertext:

\begin{equation}
ss_i
=
\operatorname{ML\mbox{-}KEM.Decaps}
(dk_i,ct_j).
\label{eq:qsv2g_kem_decaps}
\end{equation}

For a correctly generated and delivered ciphertext,
ML-KEM correctness implies matching secrets except with its
specified negligible decapsulation-failure probability
$\varepsilon_{\mathrm{KEM}}$:

\begin{equation}
\Pr[ss_i=ss_j]\ge 1-\varepsilon_{\mathrm{KEM}}.
\label{eq:qsv2g_kem_shared_secret}
\end{equation}

Key confirmation below detects a mismatched shared secret;
an invalid confirmation aborts before transaction approval.

\textit{Step 4: Session identifier and key derivation.}
Both participants independently derive a session
identifier from the authenticated transcript:

\begin{equation}
SID=
\operatorname{Trunc}_{\lambda}
\left[
H\left(
\mathrm{QSV2G\mbox{-}SID}
\parallel\mathcal{T}
\right)
\right],
\label{eq:qsv2g_session_id}
\end{equation}

where $\lambda$ specifies the identifier length.
The session identifier is public and is not
used as secret keying material.

Using the shared secret $ss$, the EV and CS
derive a pseudorandom key through HKDF:

\begin{equation}
\begin{aligned}
salt &=
H(N_E\parallel N_C),\\
PRK &=
\operatorname{HKDF\mbox{-}Extract}(salt,ss).
\end{aligned}
\label{eq:qsv2g_hkdf_extract}
\end{equation}

Two direction-specific AES-256-GCM traffic keys are derived
from the same extracted pseudorandom key, with distinct labels:

\begin{equation}
\begin{aligned}
I_E &=\mathrm{QSV2G\mbox{-}E2C}\parallel H(\mathcal T),\\
I_C &=\mathrm{QSV2G\mbox{-}C2E}\parallel H(\mathcal T),\\
K_{E\to C}&=\operatorname{HKDF\mbox{-}Expand}(PRK,I_E,32),\\
K_{C\to E}&=\operatorname{HKDF\mbox{-}Expand}(PRK,I_C,32).
\end{aligned}
\label{eq:qsv2g_session_key}
\end{equation}

Each key has 32 bytes. The EV uses $K_{E\to C}$ only to
encrypt messages to the CS, and the station uses $K_{C\to E}$
only to encrypt messages to the EV. Each sender uses a separate,
monotonically advancing 96-bit GCM nonce counter for its direction;
a counter is never rolled back or reused under the same key. The
implementation must stop transmission before counter exhaustion.
AES-256-GCM is the chosen authenticated-encryption mechanism for
this protocol; derived traffic keys are never reused as HMAC keys.

Two separate confirmation keys are derived
with distinct application labels:

\begin{equation}
\begin{aligned}
K_E^{conf}
&=\operatorname{HKDF\mbox{-}Expand}
\left(PRK,\mathrm{KC\mbox{-}EV}
\parallel H(\mathcal{T}),32\right),\\
K_C^{conf}
&=\operatorname{HKDF\mbox{-}Expand}
\left(PRK,\mathrm{KC\mbox{-}CS}
\parallel H(\mathcal{T}),32\right).
\end{aligned}
\label{eq:qsv2g_confirmation_keys}
\end{equation}

\textit{Step 5: Mutual key confirmation.}
The EV first computes a confirmation tag:

\begin{equation}
\mu_E=
\operatorname{HMAC}_{K_E^{conf}}
\left(
\mathrm{EV\mbox{-}FINISHED}
\parallel SID\parallel H(\mathcal{T})
\right).
\label{eq:qsv2g_ev_finished}
\end{equation}

It sends $(SID,\mu_E)$ to $CS_j$.
The station independently verifies $\mu_E$
using its derived confirmation key. If the
verification succeeds, $CS_j$ computes

\begin{equation}
\mu_C=
\operatorname{HMAC}_{K_C^{conf}}
\left(
\mathrm{CS\mbox{-}FINISHED}
\parallel SID\parallel H(\mathcal{T})
\parallel\mu_E
\right).
\label{eq:qsv2g_cs_finished}
\end{equation}

The station sends $\mu_C$ to the EV, which
verifies the tag using its locally derived
$K_C^{conf}$. An invalid confirmation causes
the session to be rejected, and neither
participant proceeds with authorization.

Upon successful key confirmation, both
participants retain $SID$ and the two direction-specific
traffic keys $K_{E\to C}$ and $K_{C\to E}$.
The temporary decapsulation key $dk_i$,
intermediate shared secret $ss$, and
temporary confirmation-key material are
securely erased when no longer required.
The authenticated session provides the
communication context for the subsequent
distributed authorization procedure; it
does not itself grant permission to perform
a charging or discharging operation.

\subsection{Distributed Authorization}

% Short opening paragraph
Once the session has been authenticated and its traffic keys confirmed,
$CS_j$ initiates authorization for the operation requested
by $EV_i$. The station first constructs an authorization
context that identifies the current session and specifies
the permitted transaction parameters. UC-A, UC-B, and UC-C
then evaluate the request according to their respective
authorization policies. Authorities that approve the
request participate in the threshold procedure using the
quantum shares prepared during system initialization.
A valid authorization requires contributions from at
least two distinct authorities. The resulting decision
is bound to the transaction context through an
authorization proof, which is subsequently delivered
to the EV for verification. The following steps describe
context generation, authority-side validation,
threshold reconstruction, and proof generation.

% 1. Authorization Context Generation
\noindent\textbf{1) Authorization Context Generation:}
After establishing the authenticated session, $CS_j$
determines the transaction parameters associated with
the operation requested by $EV_i$. These include the
permitted energy or power limit $LIM$ and the validity
interval $VLD$. The values must be agreed upon before
the charging station submits the request for distributed
authorization.

\textit{Step 1: Transaction parameter agreement.}
The station prepares a transaction record

\begin{equation}
\begin{aligned}
M_{\mathrm{ctx}}=\operatorname{Encode}\big(
&PID_i,CSID_j,SID,OP,N_E,N_C,\\
&TS_E,TS_C,LIM,VLD
\big),
\end{aligned}
\label{eq:qsv2g_context_record}
\end{equation}

where $TS_C$ is the station timestamp authenticated
during session establishment. The parameter $LIM$
specifies a numerical bound together with its unit,
whereas $VLD$ identifies the interval within which
the authorization may be exercised.

The station transmits $M_{\mathrm{ctx}}$ to the EV
over the authenticated session using

\begin{equation}
(C_{\mathrm{ctx}},Tag_{\mathrm{ctx}})
=
\operatorname{AEAD.Enc}_{K_{C\to E}}
(n_{\mathrm{ctx}},M_{\mathrm{ctx}},AD_{\mathrm{ctx}}),
\label{eq:qsv2g_context_offer}
\end{equation}

where $n_{\mathrm{ctx}}$ is the next 96-bit CS-to-EV nonce,
and $AD_{\mathrm{ctx}}=\operatorname{Encode}(
\mathrm{QSV2G\mbox{-}CTX},SID,CSID_j)$ is authenticated
associated data. The nonce, ciphertext, and GCM tag are delivered
together; the EV rejects the record on authentication failure.

Upon successful decryption, $EV_i$ checks the
session identifier, station identity, requested
operation, freshness values, and proposed limits.
The vehicle rejects parameters that do not match
its outstanding request or exceed its accepted
operating constraints.

If the proposed transaction is acceptable,
the EV computes $AC$ from the validated record
and generates the consent tag $\tau_E^{AC}$
as defined in (\ref{eq:qsv2g_ev_authz_consent}).
It returns an authenticated acknowledgment
containing $(SID,AC,\tau_E^{AC})$, protected under
$K_{E\to C}$ with a fresh EV-to-CS GCM nonce.
The CS verifies the acknowledgment and forwards
$\tau_E^{AC}$ unchanged with the authority request;
the station itself cannot verify the EV's registration-key tag.

\textit{Step 2: Authorization context computation.}
The station computes a transaction-specific digest
from the agreed record:

\begin{equation}
AC=
H\left(
\mathrm{QSV2G\mbox{-}AC}
\parallel M_{\mathrm{ctx}}
\right).
\label{eq:qsv2g_context}
\end{equation}

The digest covers every field in $M_{\mathrm{ctx}}$;
changing any typed field yields a different hash input.
The $\operatorname{Encode}$ function uses a
deterministic, unambiguous representation of
each field, including explicit lengths, data
types, and units. Consequently, two records
containing different transaction parameters
cannot be interpreted as the same encoded input.

The resulting digest identifies the particular
authorization request. A change in the station,
session, operation, freshness values, permitted
limit, or validity interval changes the encoded
record and, except with negligible hash-collision
probability, produces a different $AC$.

\textit{Step 3: Distribution of the authorization request.}
The charging station submits the context record
and its digest to UC-A, UC-B, and UC-C through
authenticated authority communication channels:

\begin{equation}
Req_{\mathrm{auth}}
=
(M_{\mathrm{ctx}},AC,\tau_E^{AC}).
\label{eq:qsv2g_authz_request}
\end{equation}

Each authority independently recomputes
$H(\mathrm{QSV2G\mbox{-}AC}\parallel M_{\mathrm{ctx}})$
and compares the result with the received $AC$.
A mismatch causes the request to be rejected
before policy evaluation.

The digest serves as the common reference for
the subsequent authority decisions and quantum
threshold operations. It does not, by itself,
constitute authorization evidence; an approving
authority must separately authenticate and bind
its contribution to the validated context.

% 2. Distributed Authority Verification
\noindent\textbf{2) Distributed Authority Verification:}
Upon receiving $Req_{\mathrm{auth}}$, the utility authorities
independently examine the transaction record before
contributing to the threshold authorization procedure.
Each authority evaluates the same authorization context
$AC$, but applies the policies associated with its
administrative role. The verification process is
described below.

\textit{Step 1: Verification of transaction evidence.}
UC-A retrieves the enrollment record associated with
$PID_i$ and confirms the EV authentication result
established earlier. UC-B verifies the corresponding
station authentication result, while UC-C obtains
authenticated confirmation of both results through
the inter-authority communication channels.

Each authority checks the received context and rejects expired,
unknown, or revoked state; UC-B and UC-C require an authenticated,
current consent-status statement from UC-A. The authorities compare
the transaction identifiers, operation, and freshness parameters in
$M_{\mathrm{ctx}}$ with the completed authentication
exchange. In particular, $SID$, $PID_i$, $CSID_j$,
$OP$, $N_E$, and $N_C$ must correspond to the
same authenticated session.

To provide UC-A with verifiable evidence that the EV
accepted the proposed transaction parameters,
$EV_i$ computes

\begin{equation}
\tau_E^{AC} =
\operatorname{HMAC}_{K_i^{reg}}
\left(
\mathrm{AUTHZ\mbox{-}CONSENT}
\parallel AC
\right).
\label{eq:qsv2g_ev_authz_consent}
\end{equation}

The tag is generated after the EV validates
$M_{\mathrm{ctx}}$ and is included in its
authenticated acknowledgment to $CS_j$.
The submitted $Req_{\mathrm{auth}}$ includes
$\tau_E^{AC}$, which UC-A verifies using its enrollment record. UC-A recomputes the expected
value using the registration secret associated with
$PID_i$ and checks

\begin{equation}
\begin{aligned}
\tau_{E,*}^{AC} &=\operatorname{HMAC}_{K_i^{reg}}\big(\\
&\quad \mathrm{AUTHZ\mbox{-}CONSENT}\parallel AC\big),\\
\operatorname{Consent}_{E}&=
[\tau_E^{AC}\stackrel{?}{=}\tau_{E,*}^{AC}].
\end{aligned}
\label{eq:qsv2g_ev_consent_verify}
\end{equation}

UC-A communicates its authenticated verification
result to UC-B and UC-C. Such confirmation is an attestation by
UC-A, not independently verifiable EV consent at UC-B or UC-C;
those authorities require an unrevoked and authenticated UC-A
trust relationship. A missing or invalid
consent confirmation prevents an authority from
treating the CS-submitted record as evidence of
EV agreement.

\textit{Step 2: Local policy and revocation checks.}
Let $\mathcal{P}=\{A,B,C\}$ denote the utility
authorities. Each UC-$a$ evaluates the requested
operation under its local policy
$\Pi_a(M_{\mathrm{ctx}})$.

UC-A checks the EV's registration status and
home-domain permissions. UC-B examines the
visited station's authorization to provide the
requested service, including applicable
operational limits. UC-C evaluates the
independent grid or trust constraints associated
with the transaction. The policies may differ
for charging and discharging operations.

Each authority also checks whether the EV's
pseudonym is revoked according to the revocation
information available to that authority:

\begin{equation}
\operatorname{Revoked}_a(PID_i)=0,
\qquad a\in\mathcal{P}.
\label{eq:qsv2g_revocation}
\end{equation}

The result is incorporated into the authority's
local decision. Revocation evidence from another domain must be
authenticated and sufficiently fresh under an explicitly stated
update policy; a local cache alone does not establish universal
revocation freshness. A successful authentication
exchange does not override a subsequent policy
restriction or revocation condition.

\textit{Step 3: Authority decision computation.}
For each UC-$a$, define the transaction-validity
condition

\begin{equation}
\begin{aligned}
V_a(AC)={}&
\operatorname{ContextValid}_a(AC)\\
&\land\operatorname{AuthValid}_a(SID)\\
&\land\operatorname{ConsentValid}_a(AC)\\
&\land\operatorname{TimeValid}_a(VLD),
\end{aligned}
\label{eq:qsv2g_authority_validation}
\end{equation}

where the predicates represent the authority's
verification of the submitted context, authenticated
session, EV consent evidence, and transaction
validity interval. For UC-B and UC-C, consent
validation uses the authenticated verification
result issued by UC-A.

The authority's decision is then

\begin{equation}
d_a(AC)=
\begin{cases}
1, &
V_a(AC)=1
\land \Pi_a(M_{\mathrm{ctx}})=1\\
&\land
\operatorname{Revoked}_a(PID_i)=0,\\[2pt]
0, & \text{otherwise}.
\end{cases}
\label{eq:qsv2g_authority_decision}
\end{equation}

An authority sets $d_a(AC)=1$ only when all
applicable checks succeed. A rejected, expired,
or inconsistent request yields $d_a(AC)=0$.

\textit{Step 4: Context-bound approval recording.}
An approving authority records its decision
together with the authorization context and
the identifier $\nu$ of the allocated QSS
resource:

\begin{equation}
\alpha_a^{(\nu)}
=
(a,AC,\nu,d_a(AC)).
\label{eq:qsv2g_authority_approval}
\end{equation}

Before proceeding, the participating authorities
must agree on the same unconsumed resource
identifier $\nu$. Approval records are exchanged
only through authenticated authority channels
and are associated with the pending transaction.

Each authority may issue at most one effective
approval for a given $(AC,\nu)$ pair. Duplicate
records, mismatched contexts, and approvals
associated with consumed resources are not
eligible for threshold processing.

Only authorities with $d_a(AC)=1$ proceed
to the quantum-share reconstruction procedure.
A rejecting or unavailable authority contributes
no share. This threshold rule applies after prior authentication
and authenticated EV consent verification have succeeded. In this
specified online flow UC-A and UC-B are therefore required for
pre-authorization checks even if either authority's QSS share is
not selected. The approval record $\alpha_a^{(\nu)}$
is an authenticated policy decision; it is not
itself a quantum share or an independently
EV-verifiable authorization token.

% 3. Quantum Share Reconstruction and Threshold Verification
\noindent\textbf{3) Quantum Share Reconstruction and
Threshold Verification:}
Following authority-side validation, the approved
participants execute the threshold recovery procedure
using the quantum resources prepared during system
initialization. For each authorization resource $\nu$,
the registered shares are associated with the three
utility authorities and can be recovered through any
authorized pair. The procedure is performed as follows.

\textit{Step 1: Selection of the authorized pair.}
Let $\mathcal{A}_{AC}$ denote the set of distinct
authorities that approved the current authorization
context:

\begin{equation}
\mathcal{A}_{AC}
=
\left\{
a\in\mathcal{P}
\mid d_a(AC)=1
\right\}.
\label{eq:qsv2g_approved_set}
\end{equation}

Threshold processing is permitted only when

\begin{equation}
|\mathcal{A}_{AC}|\geq 2.
\label{eq:qsv2g_approval_threshold}
\end{equation}

If fewer than two authorities approve the request,
the authorization procedure terminates without
reconstructing the protected secret.

Otherwise, two distinct authorities $u$ and $v$
are selected such that

\begin{equation}
\{u,v\}\subseteq\mathcal{A}_{AC},
\qquad
u\neq v.
\label{eq:qsv2g_selected_pair}
\end{equation}

The selected authorities verify that their approval
records refer to the same $(AC,\nu)$ pair and that
the corresponding quantum resource remains available.
Before any measurement, the authorities perform an authenticated,
crash-durable reservation transition
$\mathrm{UNUSED}\rightarrow\mathrm{RESERVED}(AC,u,v)$.
Both share holders must acknowledge the same tuple; a conflict,
uncertain state, or communication failure causes fail-closed rejection.
The resource is dedicated to this attempt and cannot be used by a
concurrent request. Reuse after any attempted measurement is forbidden.

\textit{Step 2: Quantum-share measurement.}
Each selected authority retrieves its assigned
quantum-share bundle $\mathcal{Q}_{u}^{(\nu)}$
or $\mathcal{Q}_{v}^{(\nu)}$, as defined in
(\ref{eq:qss_share_bundle}).

For every secret contribution generated by UC-$a$,
where $a\in\mathcal{P}$, the selected authorities
measure their corresponding qutrit subsystems in
the computational basis.

The measurements produce outcomes

\begin{equation}
x_{a\rightarrow u,k}^{(\nu)},
\quad
x_{a\rightarrow v,k}^{(\nu)}
\in\mathbb{Z}_3,
\label{eq:qsv2g_qss_measurements}
\end{equation}

for $1\leq k\leq L$.

The authorities exchange the required measurement
outcomes through their authenticated and confidential
communication channel. Each measurement transcript
is associated with the authorized context $AC$,
resource identifier $\nu$, and identities of the
participating authorities.

The measurement consumes the corresponding
quantum subsystems. Their outcomes must not be
accepted from another transaction or reused with
a different authorization context.

\textit{Step 3: Reconstruction of secret contributions.}
Define the authority positions as

\begin{equation}
\iota(A)=0,\qquad
\iota(B)=1,\qquad
\iota(C)=2.
\label{eq:qsv2g_authority_positions}
\end{equation}

According to the quantum-state encoding in
(\ref{eq:qss_qutrit_encoding}), the measurement
outcome of subsystem $b$ satisfies

\begin{equation}
x_{a\rightarrow b,k}^{(\nu)}
=
\left(
r_{a,k}^{(\nu)}
+
\iota(b)s_{a,k}^{(\nu)}
\right)\bmod 3.
\label{eq:qsv2g_qss_outcome_relation}
\end{equation}

Subtracting the outcomes obtained by the selected
authorities eliminates the random encoding value
$r_{a,k}^{(\nu)}$.

The reconstructed contribution is therefore

\begin{equation}
\begin{aligned}
\widehat{s}_{a,k}^{(\nu)}
={}&
\left(
x_{a\rightarrow v,k}^{(\nu)}
-
x_{a\rightarrow u,k}^{(\nu)}
\right)\\
&\cdot
\left(
\iota(v)-\iota(u)
\right)^{-1}
\bmod 3,
\end{aligned}
\label{eq:qsv2g_qss_recover_symbol}
\end{equation}

where the inverse is taken in $\mathbb{Z}_3$.
Because $u\neq v$, the denominator is nonzero
and has a multiplicative inverse.

This expression applies to each of the permitted
authority pairs and is consistent with the
pairwise recovery relations established during
QSS initialization.

For each original contributing authority,
the recovered symbols form the vector

\begin{equation}
\widehat{\mathbf{s}}_a^{(\nu)}
=
\left(
\widehat{s}_{a,1}^{(\nu)},
\ldots,
\widehat{s}_{a,L}^{(\nu)}
\right).
\label{eq:qsv2g_recovered_contribution}
\end{equation}

The combined authorization secret is recovered
componentwise as

\begin{equation}
\widehat{\mathbf{S}}_{auth}^{(\nu)}
=
\left(
\sum_{a\in\mathcal{P}}
\widehat{\mathbf{s}}_a^{(\nu)}
\right)\bmod 3.
\label{eq:qsv2g_reconstructed_secret}
\end{equation}

For correctly prepared and distributed quantum
shares, the reconstruction equations yield

\begin{equation}
\widehat{\mathbf{S}}_{auth}^{(\nu)}
=
\mathbf{S}_{auth}^{(\nu)}.
\label{eq:qsv2g_qss_correctness}
\end{equation}

\textit{Step 4: Verification of recovered contributions.}
To check the reconstructed values against their
originally registered contributions, the authorities
use the authenticated commitments
$C_a^{(\nu)}$ established during QSS setup.

For every $a\in\mathcal{P}$, the selected pair
computes

\begin{equation}
\begin{aligned}
\widehat{C}_a^{(\nu)}
=
H\big(
&\mathrm{QSV2G\mbox{-}QSS\mbox{-}COM}
\parallel a\\
&\parallel\nu
\parallel
\operatorname{Encode}
(\widehat{\mathbf{s}}_a^{(\nu)})
\big).
\end{aligned}
\label{eq:qsv2g_reconstructed_commitment}
\end{equation}

The reconstructed contribution is accepted only
when the locally computed commitment matches
the registered commitment:

\begin{equation}
V_a^{QSS}
=
\left(
\widehat{C}_a^{(\nu)}
\stackrel{?}{=}
C_a^{(\nu)}
\right).
\label{eq:qsv2g_qss_commitment_check}
\end{equation}

All three reconstructed contributions must satisfy
their respective checks:

\begin{equation}
V_{\mathrm{rec}}
=
\bigwedge_{a\in\mathcal{P}}
V_a^{QSS}.
\label{eq:qsv2g_qss_reconstruction_valid}
\end{equation}

A mismatch causes the reconstruction attempt
to be rejected. The resource is marked consumed,
and its measured quantum subsystems are not
returned to the available resource pool.

These commitment checks provide computational
consistency between recovered classical values
and registered secret contributions. They do
not, by themselves, verify every quantum-state
operation or guarantee security against arbitrary
malicious behavior during share distribution.

\textit{Step 5: Threshold reconstruction decision.}
The reconstruction is accepted only if two
distinct authorities approved the context,
their records refer to the same resource previously reserved
while unused, and all reconstructed secret contributions
pass the commitment checks.

The corresponding decision is

\begin{equation}
D_{\mathrm{QSS}}=
\begin{cases}
1, &
|\mathcal{A}_{AC}|\geq 2
\land V_{\mathrm{rec}}=1\\
&\land
\operatorname{ResourceValid}(AC,\nu)=1,\\[2pt]
0, & \text{otherwise}.
\end{cases}
\label{eq:qsv2g_qss_threshold_decision}
\end{equation}

Here $\operatorname{ResourceValid}(AC,\nu)$ means that,
\emph{before the current measurement}, the resource was unconsumed,
reserved exclusively for $(AC,u,v)$ by both participants, and no
previous attempt used it. The corresponding state transition to
$\mathrm{CONSUMED}$ is recorded after this attempt, including on
failure. A lost or uncertain acknowledgement never permits reuse.

If $D_{\mathrm{QSS}}=1$, the selected authorities
retain the validated reconstruction result
only for the subsequent authorization-proof
generation procedure. The classical measurement
outcomes and recovered secret values remain
restricted to the authorized reconstruction
participants.

The resource $\nu$ is durably marked consumed after
its reconstruction attempt, regardless of
its outcome. A successful threshold result
establishes that the selected authorities
participated in recovering the registered
secret contributions for the allocated resource.
The separate authorization proof must still
bind their authenticated approval decisions
to the specific transaction context $AC$.

% 4. Authorization Proof Generation
\noindent\textbf{4) Authorization Proof Generation:}
After successful threshold reconstruction, the participating
authorities generate a transaction-specific authorization
proof. The reconstructed secret remains within the
authorized authority pair; neither the charging station
nor the EV receives the quantum shares or the recovered
secret. Instead, the authorities issue independently
verifiable attestations binding their approval decisions
and the QSS reconstruction result to the current
authorization context.

\textit{Step 1: Reconstruction-result binding.}
Let $u,v\in\mathcal{A}_{AC}$ denote the two distinct
authorities selected for reconstruction, with
$D_{\mathrm{QSS}}=1$. The participating authorities
compute a digest of the recovered authorization secret:

\begin{equation}
\begin{aligned}
h_{\mathrm{rec}}
=
H\big(
&\mathrm{QSV2G\mbox{-}REC}
\parallel AC\parallel\nu\\
&\parallel
\operatorname{Encode}
(\widehat{\mathbf{S}}_{auth}^{(\nu)})
\big).
\end{aligned}
\label{eq:qsv2g_reconstruction_digest}
\end{equation}

The digest cryptographically binds the locally recovered secret to
the context under the assumed entropy of the original contributions.
It is an authority statement, not a publicly checkable QSS proof.
Low-entropy test vectors must not be used for real authorization
since public hashes/commitments would permit offline guessing.
The authorities verify that they obtained the same
result and that the contribution checks defined in
(\ref{eq:qsv2g_qss_reconstruction_valid}) succeeded.
The recovered secret is not included in the
authorization messages, and the recovered plaintext and measurement
outcomes are deleted after token generation and necessary audit checks.

\textit{Step 2: Authorization statement construction.}
The two authorities form a common authorization
statement containing the session identifier,
authorization context, QSS resource identifier,
participating authorities, reconstruction digest,
and validity interval:

\begin{equation}
\begin{aligned}
M_{\mathrm{auth}}
=
\operatorname{Encode}\big(
&\mathrm{QSV2G\mbox{-}AUTHZ},AC,SID,\\
&\nu,u,v,h_{\mathrm{rec}},VLD,
\mathrm{ACCEPT}
\big).
\end{aligned}
\label{eq:qsv2g_authorization_statement}
\end{equation}

The operation $OP$, permitted limit $LIM$,
participant identifiers, and session nonces
are already bound to $AC$ through
(\ref{eq:qsv2g_context}).

Each selected authority independently checks
that $M_{\mathrm{auth}}$ corresponds to its
previously accepted approval record
$\alpha_a^{(\nu)}$, that $D_{\mathrm{QSS}}=1$,
and that the selected resource has not been
assigned to another transaction.

\textit{Step 3: Independent authority attestations.}
Let $(sk_a^{sig},pk_a^{sig})$ denote the
authenticated ML-DSA signing and verification
key pair belonging to UC-$a$.

Each participating authority signs the common
authorization statement:

\begin{equation}
\sigma_a
=
\operatorname{ML\mbox{-}DSA.Sign}
\left(
sk_a^{sig},M_{\mathrm{auth}}
\right),
\quad a\in\{u,v\}.
\label{eq:qsv2g_authority_signature}
\end{equation}

An authority produces $\sigma_a$ only after
its local authorization decision is positive
and the reconstruction checks succeed.
The two signatures must originate from
distinct registered utility authorities.

The resulting authorization proof is

\begin{equation}
\begin{aligned}
AuthToken=
(&M_{\mathrm{auth}},
\sigma_u,\sigma_v).
\end{aligned}
\label{eq:qsv2g_auth_token}
\end{equation}

The token identifies the participating
authorities but contains neither their
quantum-share registers nor the recovered
authorization secret. Signatures attest the signers' reconstruction
and approval decisions; they do not constitute a publicly
verifiable proof of the quantum measurements.

\textit{Step 4: Proof collection and threshold validation.}
UC-B acts as the authorization-proof collector in the specified
online design. Its availability is required; changing collectors
would require an additional authenticated collector-election and
routing procedure, which is not defined here.
The selected authorities transmit their signed
statements to UC-B through authenticated
communication channels.

UC-B checks that both signatures refer to
the same authorization statement and verifies

\begin{equation}
\begin{aligned}
V_{\mathrm{sig}}
={}&
\operatorname{ML\mbox{-}DSA.Verify}
(pk_u^{sig},M_{\mathrm{auth}},\sigma_u)\\
&\land
\operatorname{ML\mbox{-}DSA.Verify}
(pk_v^{sig},M_{\mathrm{auth}},\sigma_v).
\end{aligned}
\label{eq:qsv2g_authorization_signatures_verify}
\end{equation}

The final authorization decision is

\begin{equation}
D_{\mathrm{auth}}=
\begin{cases}
1, &
D_{\mathrm{QSS}}=1
\land V_{\mathrm{sig}}=1\\
&\land u\neq v
\land \operatorname{Valid}(AC,\nu,VLD)=1,\\[2pt]
0, & \text{otherwise}.
\end{cases}
\label{eq:qsv2g_final_authorization_decision}
\end{equation}

Here, $\operatorname{Valid}(AC,\nu,VLD)$
requires the authorization context to match
the pending transaction, the resource to
be assigned to that transaction, and the
authorization validity interval to remain
unexpired.

UC-B forwards $AuthToken$ and the authorization
decision to $CS_j$ only if the required
verification checks succeed. Otherwise,
the request is rejected and the reserved
QSS resource is not made available for reuse.

\textit{Step 5: Proof retention and delivery.}
The authorization result is recorded against
$(SID,AC,\nu)$ so that subsequent requests
cannot claim the same consumed authorization
resource.

The charging station retains the received
$AuthToken$ with the pending session state.
It subsequently delivers the token to $EV_i$
inside the protected authorization response.

The EV independently verifies the authority
signatures using their authenticated public
verification keys and checks that the
authorization statement matches its locally
accepted transaction context. These
EV-side checks and the conditions for
executing the authorized operation are
specified in the following subsection.

\subsection{Protected Response and Execution}
\label{subsec:qsv2g_protected_execution}

% Short opening paragraph
Following the distributed authorization procedure,
$CS_j$ prepares the authorization response for $EV_i$.
The response carries the threshold authorization proof
and the transaction parameters bound to the established
session. It is protected using the previously derived
CS-to-EV traffic key $K_{C\to E}$ to prevent unauthorized
modification or disclosure during transmission.
Upon reception, the EV verifies the message integrity,
authority signatures, and agreement with its accepted
authorization context before allowing the requested
operation. The subsequent exchange is restricted by
the approved limits and validity interval, while
rejected, expired, or completed transactions are
prevented from being executed again. The following
steps specify response protection, EV-side verification,
energy exchange, and session termination.

% 1. Protected Authorization Response Generation
\noindent\textbf{1) Protected Authorization Response Generation:}
Upon receiving the authorization result from UC-B,
$CS_j$ verifies that the decision corresponds to its
pending session and authorization context $AC$.
For an accepted request, the station checks the two
authority signatures contained in $AuthToken$ using
(\ref{eq:qsv2g_authorization_signatures_verify}).
It also confirms that the authorization statement
specifies the expected $SID$, $AC$, and validity
interval $VLD$. The station proceeds with response
generation only after these checks succeed.

\textit{Step 1: Response message construction.}
The station prepares an authorization response
containing the transaction identifiers, authorization
decision, and permitted operating parameters:

\begin{equation}
\begin{aligned}
M_{\mathrm{resp}} =
\operatorname{Encode}\big(
&SID,PID_i,CSID_j,AC,OP,\\
&N_E,N_C,LIM,VLD,\\
&D_{\mathrm{auth}},AuthToken
\big).
\end{aligned}
\label{eq:qsv2g_response_message}
\end{equation}

For an accepted transaction, $D_{\mathrm{auth}}=1$
and $AuthToken$ contains the two authority
attestations defined in (\ref{eq:qsv2g_auth_token}).
For a rejected transaction, $D_{\mathrm{auth}}=0$
and the token field is set to $\bot$, indicating
that no valid authorization proof is available.
The inclusion of $AC$ allows the EV to associate
the returned decision with its previously
accepted transaction parameters.

\textit{Step 2: Authenticated response encryption.}
The station protects $M_{\mathrm{resp}}$ using
the established CS-to-EV traffic key $K_{C\to E}$ and
an authenticated encryption with associated
data (AEAD) algorithm. It selects a fresh
nonce $n_{\mathrm{resp}}$ that has not been
used previously with the same session key.
The associated data is constructed as

\begin{equation}
AD_{\mathrm{resp}} =
\operatorname{Encode}
\left(
\mathrm{QSV2G\mbox{-}RESP},
SID,CSID_j
\right).
\label{eq:qsv2g_response_ad}
\end{equation}

The protected response is generated by

\begin{equation}
\begin{aligned}
(C_{\mathrm{resp}},Tag_{\mathrm{resp}})
=
\operatorname{AEAD.Enc}_{K_{C\to E}}\big(
&n_{\mathrm{resp}},\\
&M_{\mathrm{resp}},
AD_{\mathrm{resp}}
\big).
\end{aligned}
\label{eq:qsv2g_response_encryption}
\end{equation}

Here, $C_{\mathrm{resp}}$ denotes the encrypted
authorization message, while
$Tag_{\mathrm{resp}}$ provides authentication
and integrity protection for the ciphertext
and associated data. The nonce is transmitted
in cleartext and must remain unique across
all encryption operations performed under
$K_{C\to E}$. The CS-to-EV nonce counter is advanced across both context
proposals and response/termination messages; each encryption
uses a new value. The EV-to-CS direction has its own key and
counter. All AEAD associated data are encoded canonically.

\textit{Step 3: Protected response transmission.}
The charging station transmits

\begin{equation}
Resp_{CS} =
\left(
SID,n_{\mathrm{resp}},
C_{\mathrm{resp}},
Tag_{\mathrm{resp}}
\right)
\label{eq:qsv2g_protected_response}
\end{equation}

to $EV_i$. The session identifier enables
the vehicle to select the corresponding
session key and reconstruct
$AD_{\mathrm{resp}}$ for authenticated
decryption. The station retains the
authorization result and its association
with $(SID,AC)$ in the pending transaction
record.

The transmitted response does not disclose
the reconstructed QSS secret or individual
quantum shares. The EV obtains only the
protected authorization result and, when
the operation is approved, the signed
threshold authorization proof. Verification
of the received response is performed by
the EV before any energy-transfer operation
is permitted.


% 2. EV-Side Authorization Verification
\noindent\textbf{2) EV-Side Authorization Verification:}
Upon receiving $Resp_{CS}$, the electric vehicle
$EV_i$ verifies the protected authorization result
before accepting the requested V2G operation.
The verification procedure is performed as follows.

\textit{Step 1: Authenticated response decryption.}
The EV retrieves $K_{C\to E}$ using the received
session identifier $SID$ and reconstructs the
associated data $AD_{\mathrm{resp}}$ according to
(\ref{eq:qsv2g_response_ad}). It then performs
authenticated decryption:

\begin{equation}
\begin{aligned}
M_{\mathrm{resp}}^{*}
=
\operatorname{AEAD.Dec}_{K_{C\to E}}\big(
&n_{\mathrm{resp}},\\
&C_{\mathrm{resp}},
Tag_{\mathrm{resp}},\\
&AD_{\mathrm{resp}}
\big).
\end{aligned}
\label{eq:qsv2g_response_decryption}
\end{equation}

If authentication fails, the decryption algorithm
returns $\bot$, and the response is rejected.
The EV also rejects a response nonce previously
accepted under the same session key. Only an
authenticated plaintext is processed further.

\textit{Step 2: Session and context verification.}
After decoding $M_{\mathrm{resp}}^{*}$, the EV
compares the received $SID$, $PID_i$, $CSID_j$,
$OP$, $N_E$, and $N_C$ with its locally retained
session information. The returned operating
limit $LIM$ and validity interval $VLD$ must
also agree with the transaction parameters
accepted earlier.

Using its retained transaction record
$M_{\mathrm{ctx}}$, the EV independently
recomputes the authorization context:

\begin{equation}
AC_{EV}
=
H\left(
\mathrm{QSV2G\mbox{-}AC}
\parallel M_{\mathrm{ctx}}
\right).
\label{eq:qsv2g_ev_context_verify}
\end{equation}

The received context is accepted only when

\begin{equation}
AC_{EV}\stackrel{?}{=}AC.
\label{eq:qsv2g_ev_context_match}
\end{equation}

This check ensures that the returned decision
refers to the transaction parameters previously
accepted by the EV. The vehicle additionally
verifies that its current time falls within
$VLD$, allowing for the specified clock
synchronization tolerance.

\textit{Step 3: Threshold authorization proof verification.}
For an accepted authorization response, $EV_i$
extracts $AuthToken$ and decodes the signed
statement $M_{\mathrm{auth}}$ defined in
(\ref{eq:qsv2g_authorization_statement}).

The EV first checks that the statement contains
the expected $AC$, $SID$, and $VLD$, together
with an authorization decision of
$\mathrm{ACCEPT}$. The resource identifier
$\nu$ must be associated with the current
transaction, while the two signing authorities
$u$ and $v$ must be distinct members of
$\mathcal{P}$.

Using the authenticated public verification
keys of the corresponding authorities,
the EV verifies both signatures according
to (\ref{eq:qsv2g_authorization_signatures_verify}).
The verification result is expressed as

\begin{equation}
\begin{aligned}
V_{\mathrm{token}}
={}&V_{\mathrm{sig}}
\land (u\neq v)\\
&\land \operatorname{AuthorityValid}(u,v)\\
&\land \operatorname{StatementMatch}
(M_{\mathrm{auth}},AC_{EV},SID,VLD).
\end{aligned}
\label{eq:qsv2g_ev_token_verify}
\end{equation}

Here, $\operatorname{AuthorityValid}(u,v)$
requires both verification keys to belong to
recognized utility authorities. The predicate
$\operatorname{StatementMatch}$ verifies the
agreement of the signed authorization fields
with the EV's authenticated transaction state.

The reconstruction digest $h_{\mathrm{rec}}$
is covered by the authority signatures.
The EV does not recover the QSS secret or
independently recompute this digest. Instead,
it verifies the authorities' signed
attestation to the threshold result.

\textit{Step 4: Final authorization acceptance.}
The EV accepts the authorization only when
the protected response is authentic, the
session and context checks succeed, the
authorization token is valid, and the
transaction remains eligible for execution.

Let $V_{\mathrm{sess}}$ denote the successful
session-field checks, $V_{\mathrm{ctx}}$
the agreement of the transaction parameters,
and $V_{\mathrm{time}}$ the validity-interval
check. The final EV decision is

\begin{equation}
\begin{aligned}
D_{EV}
={}&
\operatorname{AEADValid}(Resp_{CS})
\land V_{\mathrm{sess}}\\
&\land V_{\mathrm{ctx}}
\land V_{\mathrm{time}}\\
&\land V_{\mathrm{token}}
\land (D_{\mathrm{auth}}=1)\\
&\land
\operatorname{Unused}(SID,AC,\nu).
\end{aligned}
\label{eq:qsv2g_ev_final_decision}
\end{equation}

The predicate $\operatorname{Unused}$ checks that no earlier
execution was started for $(SID,AC,\nu)$ and that this transaction
is not already in a terminal state. Upon successful verification,
the EV atomically records its authorization as pending execution.
A repeated response does not issue a second execution grant.

If any verification condition fails,
$D_{EV}=0$, and the EV refuses the
requested energy exchange. A response
carrying $D_{\mathrm{auth}}=0$ is treated
as an authenticated rejection and does not
require an authorization token.

A positive $D_{EV}$ establishes the EV's
acceptance of the authority-issued permission.
Actual energy transfer remains subject to
the operating limits and execution controls
specified in the following procedure.


% 3. Authorized Energy Exchange
\noindent\textbf{3) Authorized Energy Exchange:}
After verifying the authorization response, $EV_i$
initiates the approved V2G operation through an
authenticated execution request. The charging station
permits energy transfer only when the request matches
its locally verified authorization record. Both
participants enforce the operation, permitted limits,
and validity interval throughout the transaction.
The execution procedure is described as follows.

\textit{Step 1: Execution request confirmation.}
If $D_{EV}=1$, the EV generates a fresh execution
nonce $N_X$ and constructs

\begin{equation}
M_{\mathrm{start}} =
\operatorname{Encode}
\left(
\mathrm{QSV2G\mbox{-}START},
SID,AC,\nu,OP,N_X
\right).
\label{eq:qsv2g_execution_request}
\end{equation}

The message is transmitted to $CS_j$ using the
authenticated session channel established earlier.
Its confidentiality and integrity are protected
with $K_{E\to C}$, using a fresh AEAD nonce
according to the session nonce-allocation rules.

On receiving the request, $CS_j$ verifies its
message authentication and compares $SID$, $AC$,
$\nu$, and $OP$ with the authorization statement
retained for the pending transaction. The station
also confirms that the authorization has not
expired or been consumed by another execution.

The station's execution permission is evaluated over its
locally authenticated start request and durable transaction state
and is defined as

\begin{equation}
\begin{aligned}
D_{CS}^{exec}={}&
\operatorname{ValidStart}(M_{\mathrm{start}})\\
&\land(D_{\mathrm{auth}}=1)\\
&\land\operatorname{Unused}(SID,AC,\nu)\\
&\land\operatorname{TimeValid}(VLD).
\end{aligned}
\label{eq:qsv2g_cs_execution_decision}
\end{equation}

A positive result permits the station to
activate the corresponding charging or
discharging controller. Before activation,
the transaction state is atomically changed
from \textsc{authorized} to \textsc{active},
preventing a second request from starting
the same authorized operation.

\textit{Step 2: Operation and limit enforcement.}
Let the authorization limit be represented by

\begin{equation}
LIM=(L_{\mathrm{type}},L_{\max},U_L),
\label{eq:qsv2g_execution_limit}
\end{equation}

where $L_{\mathrm{type}}$ identifies the
quantity being restricted, $L_{\max}$
specifies its maximum authorized value,
and $U_L$ denotes the corresponding unit.

For an energy-limited transaction, let
$P_{\mathrm{tr}}(t)\geq 0$ denote the
measured magnitude of power transferred
between the EV and station. The cumulative
transferred energy is

\begin{equation}
E_{\mathrm{tr}}(t)
=
\int_{t_0}^{t}
P_{\mathrm{tr}}(\tau)\,d\tau,
\label{eq:qsv2g_transferred_energy}
\end{equation}

where $t_0$ is the execution start time.
The station must enforce

\begin{equation}
E_{\mathrm{tr}}(t)\leq E_{\max},
\qquad
E_{\max}=L_{\max},
\label{eq:qsv2g_energy_constraint}
\end{equation}

when $L_{\mathrm{type}}=\mathrm{ENERGY}$.

For a power-limited authorization,
the corresponding condition is

\begin{equation}
P_{\mathrm{tr}}(t)\leq P_{\max},
\qquad
P_{\max}=L_{\max},
\label{eq:qsv2g_power_constraint}
\end{equation}

when $L_{\mathrm{type}}=\mathrm{POWER}$.
The appropriate constraint is selected
from the signed authorization context.
Any additional equipment or grid safety
limits remain independently enforceable.

The operation $OP$ determines the permitted
direction of transfer. For
$OP=\mathrm{CHARGE}$, energy flows from
the charging infrastructure to the EV.
For $OP=\mathrm{DISCHARGE}$, the EV
supplies energy to the grid through the
authorized station interface. A change
of direction requires a new authorization.

\textit{Step 3: Validity and execution monitoring.}
Let the authorized validity interval be

\begin{equation}
VLD=[t_{\mathrm{begin}},t_{\mathrm{end}}],
\label{eq:qsv2g_execution_validity}
\end{equation}

with $t_{\mathrm{begin}}<t_{\mathrm{end}}$.
Energy exchange is permitted only while
the execution time lies within this interval.

The station continuously evaluates the
execution condition

\begin{equation}
\begin{aligned}
\operatorname{Permit}(t)
={}&
\operatorname{Active}(SID,AC,\nu)\\
&\land
(t_{\mathrm{begin}}\leq t\leq t_{\mathrm{end}})\\
&\land
\operatorname{DirectionValid}(OP,t)\\
&\land
\operatorname{LimitValid}(LIM,t)\\
&\land
\operatorname{SafetyValid}(t).
\end{aligned}
\label{eq:qsv2g_execution_permission}
\end{equation}

The predicate $\operatorname{LimitValid}$
evaluates the applicable energy or power
constraint, while
$\operatorname{SafetyValid}$ represents
the independent operational interlocks
of the charging infrastructure.

If any condition becomes false, $CS_j$
must stop further authorized transfer
and place the transaction in a
non-active state. The EV may independently
request termination if its local operating
conditions no longer permit the exchange.

The limits must be enforced using trusted
measurements and a controller capable of
interrupting energy transfer. Cryptographic
authorization alone does not guarantee
physical compliance with these constraints.

\textit{Step 4: Execution record and replay prevention.}
During execution, the station maintains a
transaction record associated with
$(SID,AC,\nu)$, including the operation,
measured energy, and execution state.

The persistent state machine uses the following forward-only
execution transitions (a failed or expired authorization is terminal):

\begin{equation}
\begin{aligned}
\textsc{authorized}
&\rightarrow \textsc{active},\\
\textsc{active}
&\rightarrow
\{\textsc{completed},
  \textsc{interrupted}\}.
\end{aligned}
\label{eq:qsv2g_execution_states}
\end{equation}

An authorization that has entered the
\textsc{active}, \textsc{completed}, or
\textsc{interrupted} state cannot initiate
a new execution using the same
$(SID,AC,\nu)$ identifier. Repeated start
messages are rejected without restarting
the physical transfer.

The final execution outcome and associated
metering information are retained for
session termination and transaction
recording. Any subsequent charging or
discharging operation requires a new
authorization context and a fresh
authorization procedure.

% 4. Session Termination and Exceptional Handling
\noindent\textbf{4) Session Termination and Exceptional Handling:}
The V2G session terminates when the authorized energy
transfer is completed, interrupted, or no longer
permitted under the established transaction conditions.
The charging station records the final execution
state and communicates the outcome to the EV through
the authenticated session channel. Termination and
exceptional conditions are handled as follows.

\textit{Step 1: Transaction finalization.}
Let $t_f$ denote the time at which energy transfer
ends. The charging station obtains the final
metered energy $E_{\mathrm{tr}}(t_f)$ and determines
the execution status $S_f$ as

\begin{equation}
S_f\in
\{\mathrm{COMPLETED},\mathrm{INTERRUPTED}\}.
\label{eq:qsv2g_final_status}
\end{equation}

The status $\mathrm{COMPLETED}$ indicates that
the authorized transaction ended normally,
whereas $\mathrm{INTERRUPTED}$ represents
termination before normal completion because
of an operational, communication, or
authorization-related condition.

The station prepares the final transaction record

\begin{equation}
\begin{aligned}
M_{\mathrm{term}}=
\operatorname{Encode}\big(
&SID,AC,\nu,OP,S_f,\\
&t_f,E_{\mathrm{tr}}(t_f)
\big).
\end{aligned}
\label{eq:qsv2g_termination_record}
\end{equation}

The recorded energy represents the measured
quantity exchanged during the transaction.
Its accuracy depends on the integrity and
calibration of the associated metering system.

\textit{Step 2: Protected termination confirmation.}
The station generates a fresh AEAD nonce
$n_{\mathrm{term}}$ and protects the final record
using the established session key:

\begin{equation}
\begin{aligned}
(C_{\mathrm{term}},Tag_{\mathrm{term}})
=
\operatorname{AEAD.Enc}_{K_{C\to E}}\big(
&n_{\mathrm{term}},\\
&M_{\mathrm{term}},
AD_{\mathrm{term}}
\big),
\end{aligned}
\label{eq:qsv2g_termination_protection}
\end{equation}

where

\begin{equation}
AD_{\mathrm{term}}=
\operatorname{Encode}
(\mathrm{QSV2G\mbox{-}TERM},SID,CSID_j).
\label{eq:qsv2g_termination_ad}
\end{equation}

The nonce is allocated using the persistent CS-to-EV
96-bit counter and the existing $K_{C\to E}$ key.
The station transmits the protected record
to $EV_i$.

Upon receipt, the EV verifies the AEAD tag
and confirms that $SID$, $AC$, $\nu$, and
$OP$ correspond to its active transaction.
It records the authenticated termination
status and the reported metering information.
A receipt acknowledgment may be returned
over the same protected channel.

The station does not postpone a required
physical shutdown while waiting for this
acknowledgment. If receipt delivery fails,
the termination record remains available
for authenticated reconciliation.

\textit{Step 3: Expiry and interruption handling.}
An active transaction is interrupted if its
validity interval expires, the permitted
operating constraints are violated, or a
safety condition requires shutdown.
Loss of the protected communication channel
also prevents further cryptographically
authorized control commands and triggers
the configured safe-stop procedure.

The terminal execution state is recorded
against $(SID,AC,\nu)$, and the corresponding
authorization cannot initiate another
energy-transfer operation. A repeated
execution or termination message must not
reactivate a completed or interrupted
transaction.

If an interruption occurs before the
termination record can be exchanged, the
station retains the locally observed
execution outcome. The EV and station may
later reconcile their records through an
authenticated communication channel.
A new charging or discharging attempt
requires a new session authorization rather
than reuse of the earlier $AuthToken$.

\textit{Step 4: Revocation and resource closure.}
If an authority reports a relevant credential
revocation or withdrawal of authorization,
the station applies the corresponding
termination policy. A revocation received
during an active transaction must be
authenticated and checked against the
affected participant and authorization
context before the station changes the
execution state.

The consumed QSS resource $\nu$ remains
unavailable for future reconstruction.
The authorities retain the minimum records
required to associate its use with the
corresponding authorization decision.

After termination processing, the EV and
CS invalidate the session for further
execution requests, persist the terminal transaction state, and
erase $K_{E\to C}$ and $K_{C\to E}$ when no further authenticated
reconciliation is required. If reconciliation remains pending,
separate restricted retention rules must govern the keys and records. Transaction records required for
accounting, auditing, or dispute resolution
may be retained according to the applicable
retention policy.

Any exceptional request to resolve a
pseudonym to a registered EV identity is
handled by UC-A under the established
disclosure policy. Such disclosure is
separate from routine session termination
and does not expose the QSS authorization
secret or the EV--CS session key.

\section{Security Analysis} \label{sec5}
\subsection{Security Features Supported}
\subsection{Security Verification}
\subsection{Security Features Comparison}

\section{Performance Evaluation} \label{sec6}
\subsection{Simulation Setup and Benchmarks}
\subsection{Results and Analysis}
\subsection{Computation Overhead} \label{comp}
\subsection{Energy Consumption} \label{ene}
\subsection{Communication Overhead} \label{comm}

\section{Conclusion and Future work} \label{sec7}

\bibliographystyle{IEEEtran}
\bibliography{new}
\end{document}
